% ============================================================
%  ON THE CONVERGENCE OF THE ORTHOGONAL SPECTRAL SEQUENCE
%  César Galindo  –  IHES-style LaTeX  (English translation)
%  Original: Tesis doctoral, UNAM, October 2022
%  Director: Dr. José Pablo Peláez Menaldo
%            Chebyshev Laboratory, St.\ Petersburg State University / IMATE-CU
% ============================================================

\documentclass[12pt,a4paper]{article}

% ------------------------------------------------------------------
% IHES-style packages
% ------------------------------------------------------------------
\usepackage[T1]{fontenc}
\usepackage[utf8]{inputenc}
\usepackage{lmodern}
\usepackage{microtype}

% Mathematics
\usepackage{amsmath,amssymb,amsthm,mathtools}
\usepackage{stmaryrd}          % \llbracket, \rrbracket
\usepackage{mathrsfs}          % \mathscr

% Adjoint pair arrow (amssymb)
\providecommand{\rightleftarrows}{%
  \mathrel{\rightarrow\mkern-14mu\leftarrow}}

% Layout  (IHES: generous margins, single column)
\usepackage[top=3cm, bottom=3cm, left=3.2cm, right=3.2cm]{geometry}

% Cross-references & hyperlinks
\usepackage[colorlinks=true,
            linkcolor=blue!60!black,
            citecolor=green!50!black,
            urlcolor=blue!70!black]{hyperref}
\usepackage{cleveref}

% Diagrams
\usepackage{tikz-cd}
\usepackage{tikz}
\usetikzlibrary{arrows,calc,patterns,decorations.pathmorphing}

% Enumeration
\usepackage{enumitem}

% Header / footer
\usepackage{fancyhdr}
\pagestyle{fancy}
\fancyhf{}
\fancyhead[LE,RO]{\small\thepage}
\fancyhead[LO]{\small\itshape On the Convergence of the Orthogonal Spectral Sequence}
\fancyhead[RE]{\small\itshape C.\ Galindo}
\renewcommand{\headrulewidth}{0.4pt}

% Spacing
\usepackage{setspace}
\setstretch{1.08}

% ------------------------------------------------------------------
% THEOREM ENVIRONMENTS  (IHES style: definition font for statements)
% ------------------------------------------------------------------
\newtheorem{theorem}{Theorem}[section]
\newtheorem{proposition}[theorem]{Proposition}
\newtheorem{lemma}[theorem]{Lemma}
\newtheorem{corollary}[theorem]{Corollary}

\theoremstyle{definition}
\newtheorem{definition}[theorem]{Definition}
\newtheorem{example}[theorem]{Example}
\newtheorem{remark}[theorem]{Remark}
\newtheorem{note}[theorem]{Note}

% ------------------------------------------------------------------
% MACROS
% ------------------------------------------------------------------
% Categories
\newcommand{\Ab}{\mathbf{Ab}}
\newcommand{\Sm}{\mathrm{Sm}}
\newcommand{\Sch}{\mathrm{Sch}}
\newcommand{\Cor}{\mathrm{Cor}}
\newcommand{\PST}{\mathit{PST}}
\newcommand{\DM}{\mathit{DM}}
\newcommand{\DMeff}{DM^{\mathrm{eff}}_{-}}
\newcommand{\DMgm}{DM_{\mathrm{gm}}}
\newcommand{\DMeffgm}{\widehat{DM}^{\mathrm{eff}}_{\mathrm{gm}}}
\newcommand{\SH}{\mathcal{SH}}
\newcommand{\SHeff}{\mathcal{SH}^{\mathrm{eff}}}
\newcommand{\ShvNis}{\mathrm{Shv}_{\mathrm{Nis}}}
\newcommand{\PreShv}{\mathrm{PreShv}}

% Functors / operations
\newcommand{\bc}{bc}
\newcommand{\hocolim}{\operatorname*{hocolim}}
\newcommand{\colim}{\operatorname*{colim}}
\newcommand{\holim}{\operatorname*{holim}}
\newcommand{\Hom}{\operatorname{Hom}}
\newcommand{\Ext}{\operatorname{Ext}}
\newcommand{\id}{\mathrm{id}}

% Graded pieces / filtration
\newcommand{\Gr}{\mathrm{Gr}}
\newcommand{\CH}{\mathrm{CH}}

% Spectral sequence terms
\newcommand{\Einf}{E_{\infty}}

% Number fields
\newcommand{\Z}{\mathbb{Z}}
\newcommand{\Q}{\mathbb{Q}}
\newcommand{\A}{\mathbb{A}}
\newcommand{\PP}{\mathbb{P}}
\newcommand{\Gm}{\mathbb{G}_m}

% Orthogonal complement
\newcommand{\DMperp}[1]{DM^{\perp}(#1)}

% ------------------------------------------------------------------
\begin{document}

% ------------------------------------------------------------------
%  TITLE PAGE
% ------------------------------------------------------------------
\thispagestyle{empty}

\begin{center}

\vspace*{1.5cm}

{\Large\bfseries
ON THE CONVERGENCE OF THE\\[4pt]
ORTHOGONAL SPECTRAL SEQUENCE
}

\vspace{1.8cm}

{\large César Galindo}

\vspace{0.4cm}

{\normalsize
Instituto de Matemáticas, UNAM\\
Av.\ Universidad s/n, Ciudad de México\\
Chebyshev Laboratory, St.\ Petersburg State University
}

\vspace{2cm}

\begin{tabular}{rl}
\textit{Advisor:} &
Dr.\ José Pablo Peláez Menaldo \\[2pt]
& \small Chebyshev Laboratory, St.\ Petersburg State University \\[2pt]
& \small \url{https://chebyshev.spbu.ru/en/people/jose-pablo-pelaez-menaldo/} \\[8pt]

\textit{Committee:} &
Dr.\ Enrique Javier Elizondo Huerta (IMATE--CU) \\[2pt]
&
Dr.\ Felipe de Jesús Zaldívar Cruz (UAM--Iztapalapa)
\end{tabular}

\vspace{2cm}

{\normalsize
Doctoral Thesis — Universidad Nacional Autónoma de México\\
October 2022
}

\end{center}

\newpage

% ------------------------------------------------------------------
%  ABSTRACT
% ------------------------------------------------------------------
\begin{abstract}
\noindent
One of the central theorems in Algebraic Geometry is the Néron–Severi theorem, which
asserts that for a smooth projective variety $X$ over an algebraically closed field,
the group $\CH^1(X)/\CH^1_{\mathrm{num}}(X)$ is finitely generated.  Motivated by
this, Bloch, Beilinson, and Murre independently introduced a conjectural filtration
on the higher Chow groups $\CH^r(X)$, $r>1$.  To study finite filtrations on these
groups, Peláez introduced in~\cite{Pel17} a tower of triangulated functors
$bc_{\leq\bullet}$ in the Voevodsky category of mixed motives $DM$ and associated a
spectral sequence
\[
  E^1_{p,q} = \Hom_{DM}\!\bigl(B,\,(bc_{p/p-1}A)[q-p]\bigr)
  \;\Rightarrow\; \Hom_{DM}(B,A).
\]
The main object of this thesis is to study the convergence of this \emph{orthogonal
spectral sequence}.  The main result (Theorem~\ref{thm:main}) shows that this
spectral sequence converges strongly when $B = M(X)(r)[s]$, with $r,s\in\Z$ and $X$
any smooth scheme of finite type over $k$.  We also show that the analogous result
fails in the $\mathbb{A}^1$-stable homotopy category $\SH$ of Morel–Voevodsky
(Proposition~\ref{prop:counterexample}), and we identify conditions on $A\in\SH$
under which strong convergence holds (Proposition~\ref{prop:SHconvergence}).  As a
direct consequence, we obtain a spectral sequence that converges strongly at the
$E_1$-term of the Voevodsky slice spectral sequence
(Corollary~\ref{cor:sliceconv}).
\end{abstract}

\tableofcontents
\newpage

% ==================================================================
\section*{Introduction}
\addcontentsline{toc}{section}{Introduction}
% ==================================================================

One of the central theorems in Algebraic Geometry is the \emph{Néron–Severi theorem}
\cite{Sev34,Ner52a,Ner52b}: for a smooth projective variety $X$ over an algebraically
closed field, the group $\CH^1(X)/\CH^1_{\mathrm{num}}(X)$ is finitely generated,
where $\CH^1(X)$ is the group of algebraic cycles of codimension~$1$ modulo rational
equivalence and $\CH^1_{\mathrm{num}}(X)$ is the subgroup of numerically trivial
cycles.  Since $\CH^1_{\mathrm{num}}(X)$ is the group of rational points of the Picard
variety of $X$ (an abelian variety), this result decomposes $\CH^1(X)$ into a
continuous component — parametrised by an abelian variety — and a discrete component
$\CH^1(X)/\CH^1_{\mathrm{num}}(X)$ which is a finitely generated abelian group.

Motivated by this, the expectation was that for the higher Chow groups $\CH^r(X)$,
$r>1$, there would similarly exist a subgroup parametrised by some naturally associated
algebraic variety such that the quotient was finitely generated.  However, this was
dashed by Mumford~\cite{Mum69}, who proved that for a surface $S$ with
$\dim H^0(S,\Omega_S^2)>0$, the group of zero cycles $\CH^2(S)$ is ``enormous'' and
in particular admits no such parametrisation.

This led experts to conclude that the Chow groups $\CH^r(X)$ for $r>1$ are in
general chaotic.  Nevertheless, in the late 1970s Bloch, Beilinson, and Murre
independently introduced~\cite{Blo80,Bei87,Mur93} a conjectural programme in which the
algebraic cycle groups $\CH^r(X)$, $r>1$, have accessible properties.  The central
piece of this programme is the existence of a filtration on $\CH^r(X)$ whose
associated graded pieces have accessible properties.

In order to study finite filtrations on the Chow groups of a smooth projective variety
$Y$ over a field $k$ satisfying some properties of the (still conjectural)
Bloch–Beilinson–Murre filtration, Peláez introduced in~\cite{Pel14} a tower of
triangulated functors $bc_{\leq\bullet}$:
\[
  \cdots \to bc_{\leq n-1} \to bc_{\leq n} \to bc_{\leq n+1} \to \cdots
\]
in the Voevodsky triangulated category of mixed motives $DM$.  The filtration on the
Chow groups with coefficients in a commutative ring $R$ is defined by evaluating the
orthogonal tower $bc_{\leq\bullet}$ on the motive of a point $\Z_R$ and then
considering morphisms $M(Y)(-q)[-2q] \to bc_{\leq\bullet}(\Z_R)$, where $M(Y)$ is
the motive of $Y$ and $(-q)$ (resp.\ $[-2q]$) denotes the Tate twist in $DM$
(resp.\ the suspension in $DM$).  The previous process provides a filtration on the
Chow groups since
\[
  \CH^q(Y)_R \cong \Hom_{DM}\!\bigl(M(Y)(-q)[-2q],\,\Z_R\bigr).
\]

Given $A, B \in DM$, it is possible to evaluate the tower on $A$ and then map $B$
into $bc_{\leq\bullet}(A)$ to obtain a spectral sequence:
\begin{equation}\label{eq:SS}
  E^1_{p,q} = \Hom_{DM}\!\bigl(B,\,(bc_{p/p-1}A)[q-p]\bigr)
  \;\Rightarrow\; \Hom_{DM}(B,A),
\end{equation}
where $bc_{p/p-1}A$ is defined in terms of a canonical distinguished triangle in
$DM$ (\cite{Pel17}, 3.2.8):
\[
  bc_{\leq p-1}A \to bc_{\leq p}A \to bc_{p/p-1}A.
\]

The main goal of this thesis is to study convergence properties of this spectral
sequence.  The main result (Theorem~\ref{thm:main}) shows that the spectral
sequence~\eqref{eq:SS} converges strongly for $B = M(X)(r)[s]$, with $r,s\in\Z$
and $X$ any smooth scheme of finite type over $k$.

We also show that the analogous results fail in the $\mathbb{A}^1$-stable homotopy
category of Morel–Voevodsky $\SH$, providing explicit counterexamples
(Proposition~\ref{prop:counterexample} and Corollary~\ref{cor:counterexample2}).
On the other hand, we verify that, under appropriate conditions on $A\in\SH$, the
spectral sequence is strongly convergent
(Proposition~\ref{prop:SHconvergence}), and that these conditions are satisfied
for the Voevodsky slices $s_mG\in\SH$.  As a direct consequence, we obtain a
spectral sequence strongly convergent at the $E_1$ term of the Voevodsky slice
spectral sequence (Corollary~\ref{cor:sliceconv}).

% ==================================================================
\section{Preliminaries}
% ==================================================================

% ------------------------------------------------------------------
\subsection{Triangulated categories}
% ------------------------------------------------------------------

The main references for this section are~\cite{Nee96,BBD}.

Let $\mathcal{A}$ be an additive category equipped with an autoequivalence
$\Sigma:\mathcal{A}\to\mathcal{A}$ of categories.  For every integer $n\in\Z$ and
every object $X$ (resp.\ morphism $f$) of $\mathcal{A}$, we define
$X[n]:=\Sigma^n(X)$ (resp.\ $f[n]=\Sigma^n(f)$).

\begin{definition}[Verdier]\label{def:triangulated}
Let $\mathcal{A}$ be an additive category.  The structure of a \emph{triangulated
category} is given by an autoequivalence $\Sigma=[1]:\mathcal{A}\to\mathcal{A}$,
called the \emph{translation functor}, and a class of diagrams in $\mathcal{A}$.
A \emph{triangle} of $(\mathcal{A},\Sigma)$ is a sequence of morphisms
\[
  X \xrightarrow{u} Y \xrightarrow{v} Z \xrightarrow{w} X[1].
\]
A morphism between distinguished triangles is a triple $(f,g,h)$ of morphisms of
$\mathcal{A}$ making the following diagram commute:
\[
\begin{tikzcd}
  X \ar[r,"u"] \ar[d,"f"] & Y \ar[r,"v"] \ar[d,"g"] & Z \ar[r,"w"] \ar[d,"h"] & X[1] \ar[d,"f{[1]}"]\\
  X' \ar[r,"u'"] & Y' \ar[r,"v'"] & Z' \ar[r,"w'"] & X'[1].
\end{tikzcd}
\]
Such a triple $(f,g,h)$ is an isomorphism if the vertical arrows are isomorphisms.
We require the distinguished triangles to satisfy the following axioms:
\begin{enumerate}[label=\textup{(TR\arabic*)}]
  \item\label{TR1}
    \begin{enumerate}[label=(\alph*)]
      \item Every triangle of the form $X\xrightarrow{1_X} X\to 0\to X[1]$ is
        distinguished.
      \item Every triangle isomorphic to a distinguished triangle is distinguished.
      \item For every morphism $f:X\to Y$ in $\mathcal{A}$, there exists a
        distinguished triangle of the form
        $X\xrightarrow{f} Y \to Z \to X[1]$.
    \end{enumerate}
  \item\label{TR2}
    A triangle $X\xrightarrow{u} Y\xrightarrow{v} Z\xrightarrow{w} X[1]$ is
    distinguished if and only if
    $Y\xrightarrow{v} Z\xrightarrow{-w} X[1]\xrightarrow{-\Sigma u} Y[1]$
    is distinguished.
  \item\label{TR3}
    Suppose we have a commutative diagram of distinguished triangles with vertical
    arrows $f$ and $g$.  Then there exists (not necessarily unique) an arrow
    $h:Z\to Z'$ making the full diagram commute.
  \item\label{TR4} \textbf{Octahedral axiom.}
    (We refer to~\cite{BBD} for the precise statement.)
\end{enumerate}
\end{definition}

\begin{note}\label{note:opposite}
If $\mathcal{T}=(\mathcal{A},\Sigma)$ is a triangulated category, then
$\mathcal{T}^{op}=(\mathcal{A}^{op},\Sigma^{-1})$ is also a triangulated category.
\end{note}

\begin{proposition}\label{prop:coproducts}
Let $\mathcal{T}$ be a triangulated category.  Suppose $\{X_\lambda\}_{\lambda\in\Lambda}$
is a set of objects of $\mathcal{T}$ and that the coproduct
$\coprod_{\lambda\in\Lambda} X_\lambda$ exists in $\mathcal{T}$.  Then the natural map
\[
  \coprod_{\lambda\in\Lambda} X_\lambda[1]
  \longrightarrow
  \Bigl(\coprod_{\lambda\in\Lambda} X_\lambda\Bigr)[1]
\]
is an isomorphism.
\end{proposition}

\begin{definition}\label{def:homological}
Let $\mathcal{T}$ be a triangulated category and $\mathcal{A}$ an abelian category.
A functor $H:\mathcal{T}\to\mathcal{A}$ is \emph{homological} if for each
distinguished triangle $X\xrightarrow{u} Y\xrightarrow{v} Z\xrightarrow{w} X[1]$
the sequence $H(X)\xrightarrow{H(u)} H(Y)\xrightarrow{H(v)} H(Z)$ is exact in
$\mathcal{A}$.
\end{definition}

By axiom~\ref{TR2} the sequence extends to a long exact sequence in both directions:
\[
  \cdots\to H(Z[-1])\xrightarrow{H(w[-1])} H(X)\xrightarrow{H(u)} H(Y)
  \xrightarrow{H(v)} H(Z)\xrightarrow{H(w)} H(X[1])\to\cdots
\]

\begin{note}
A homological functor on $\mathcal{T}^{op}$ is a \emph{cohomological functor} on
$\mathcal{T}$.
\end{note}

\begin{lemma}\label{lem:Hom-homological}
Let $\mathcal{T}$ be a triangulated category and $U$ an object of $\mathcal{T}$.
The functor $\Hom(U,-):\mathcal{T}\to\Ab$ is homological.
\end{lemma}

\begin{proof}
Let $X\xrightarrow{u} Y\xrightarrow{v} Z\xrightarrow{w} X[1]$ be a distinguished
triangle.  We need exactness at $\Hom(U,Y)$.  First, $vu=0$: by~\ref{TR1}(a) and
\ref{TR3} applied to the canonical triangle $X\xrightarrow{\id}X\to 0\to X[1]$ and
our triangle, we obtain a morphism $h:U\to X$ (via the map $U\to 0$) with $h=0$, and
the commutative square forces $vu=0$.  For exactness, if $vf=0$ then by~\ref{TR2}
and~\ref{TR3} one constructs $h:U\to X$ with $uh=f$; see~\cite{BBD} for details.
\end{proof}

\begin{note}
It follows that $\Hom(-,U):\mathcal{T}^{op}\to\Ab$ is cohomological.
\end{note}

\begin{proposition}\label{prop:TC-properties}
Let $\mathcal{T}$ be a triangulated category and $(f,g,h)$ a morphism of triangles.
\begin{enumerate}
  \item If two of the three vertical morphisms are isomorphisms, so is the third.
  \item If $X\xrightarrow{u} Y\xrightarrow{v} Z\xrightarrow{0} X[1]$ is a
    distinguished triangle, then $u$ is an isomorphism if and only if $Z\cong 0$.
  \item Every triangle of the form $X\to Y\to Z\xrightarrow{0} X[1]$ is isomorphic
    to $X\to X\oplus Z\to Z\xrightarrow{0} X[1]$.
\end{enumerate}
\end{proposition}
\begin{proof}
See~\cite{BBD}, IV~\S1.
\end{proof}

\begin{definition}\label{def:tensor-structure}
A \emph{tensor structure} on an additive category $\mathcal{A}$ is given by a
bifunctor $\otimes:\mathcal{A}\times\mathcal{A}\to\mathcal{A}$ which is associative,
commutative and unital; that is, there is a unit object $1$ and isomorphisms
$X\otimes(Y\otimes Z)\xrightarrow{\sim}(X\otimes Y)\otimes Z$,
$X\otimes Y\xrightarrow{\sim} Y\otimes X$, $1\otimes X\xrightarrow{\sim} X$,
$X\otimes 1\xrightarrow{\sim} X$, making the obvious diagrams commute.
\end{definition}

\begin{definition}\label{def:graded-tensor}
If $\mathcal{A}$ has a translation $X\mapsto\Sigma X$, the tensor structure is
\emph{graded} if
\begin{enumerate}
  \item $\Sigma\circ(-\otimes-) = \Sigma(-)\otimes-$ and
    $\Sigma^2(-)\otimes-= -\otimes\Sigma^2(-)$ as functors
    $\mathcal{A}\times\mathcal{A}\to\mathcal{A}$.
  \item The natural isomorphisms
    $X\otimes\Sigma Y\cong\Sigma X\otimes Y=\Sigma(X\otimes Y)=\Sigma X\otimes Y$
    satisfy the condition that the composition
    $X\otimes\Sigma^2 Y\to X\otimes\Sigma^2 Y$ is the identity.
\end{enumerate}
\end{definition}

\begin{definition}\label{def:tensor-triangulated}
A \emph{tensor triangulated category} is a triangulated category
$\mathcal{T}$ which is also a tensor category (with $\otimes$ graded) such that for
each distinguished triangle
$X\xrightarrow{u} Y\xrightarrow{v} Z\xrightarrow{w} X[1]$
and every $W\in\mathcal{T}$, the sequence
\[
  X\otimes W\xrightarrow{u\otimes\id_W} Y\otimes W
  \xrightarrow{v\otimes\id_W} Z\otimes W
  \xrightarrow{w\otimes\id_W} X[1]\otimes W = (X\otimes W)[1]
\]
is a distinguished triangle.
\end{definition}

\begin{definition}\label{def:triangulated-functor}
Let $(\mathcal{T},\Sigma)$ and $(\mathcal{T}',\Sigma')$ be triangulated categories.
An additive functor $F:\mathcal{T}\to\mathcal{T}'$ is \emph{triangulated} if there
is a natural isomorphism $F\circ\Sigma\cong\Sigma'\circ F$ such that for every
distinguished triangle $X\xrightarrow{u}Y\xrightarrow{v}Z\xrightarrow{w}X[1]$
in $\mathcal{T}$, the triangle
$FX\xrightarrow{u}FY\xrightarrow{v}FZ\xrightarrow{w}F(X[1])$ is distinguished in
$\mathcal{T}'$.
\end{definition}

\begin{definition}\label{def:triangulated-subcategory}
A \emph{triangulated subcategory} of $\mathcal{T}$ is a subcategory
$\iota:\mathcal{S}\subset\mathcal{T}$ with the structure of a triangulated category
such that the inclusion functor $\iota$ is triangulated.
\end{definition}

\begin{proposition}\label{prop:adjoint-triangulated}
Let $\mathcal{T}$ and $\mathcal{T}'$ be triangulated categories, and consider an
adjoint pair of functors where $F$ is the left adjoint:
\[
\mathcal{T}\;\overset{F}{\underset{G}{\rightleftarrows}}\;\mathcal{T}'.
\]
Then $F$ is triangulated if and only if $G$ is triangulated.
\end{proposition}
\begin{proof}
See~\cite{Nee01}, Lem.\ 5.3.6.
\end{proof}

\begin{proposition}\label{prop:full-subcategory}
If $\mathcal{S}\subset\mathcal{T}$ is a full subcategory of a triangulated category
$\mathcal{T}$, then $\mathcal{S}$ is a triangulated subcategory if and only if it is
invariant under the translation functor and for every distinguished triangle
$X\to Y\to Z\to X[1]$ in $\mathcal{T}$ with $X,Y\in\mathcal{S}$, the object $Z$ is
isomorphic to an object in $\mathcal{S}$.
\end{proposition}

% ------------------------------------------------------------------
\subsubsection{Verdier localisation}

\begin{definition}\label{def:kernel}
Let $F:\mathcal{D}\to\mathcal{T}$ be a triangulated functor.  The \emph{kernel} of
$F$ is the full subcategory $\mathcal{C}=\{X\in\mathcal{D}: F(X)\cong 0\}$.
\end{definition}

\begin{proposition}\label{prop:kernel-properties}
The kernel of a triangulated functor $F:\mathcal{D}\to\mathcal{T}$ satisfies:
\begin{enumerate}
  \item $\mathcal{C}$ is a triangulated subcategory of $\mathcal{D}$.
  \item If $X\oplus Y\in\mathcal{C}$, then $X,Y\in\mathcal{C}$.
\end{enumerate}
\end{proposition}
\begin{proof}
See~\cite{Nee01}, Lem.\ 2.1.4 and Lem.\ 2.1.5.
\end{proof}

\begin{definition}\label{def:thick}
A full triangulated subcategory $\mathcal{C}$ of a triangulated category $\mathcal{D}$
is \emph{thick} if $\mathcal{C}$ is closed under direct summands.
\end{definition}

\begin{theorem}[Verdier~\cite{Ver77}]\label{thm:Verdier-localisation}
Let $\mathcal{T}$ be a small triangulated category and $\mathcal{C}\subset\mathcal{T}$
a triangulated subcategory.  There exists a triangulated category $\mathcal{T}/\mathcal{C}$
and a triangulated functor $\mathcal{Q}_\mathcal{C}:\mathcal{T}\to\mathcal{T}/\mathcal{C}$
such that, for every triangulated functor $F:\mathcal{T}\to\mathcal{T}'$ with
$\mathcal{C}$ contained in the kernel of $F$, there exists a unique triangulated
functor $\bar{F}:\mathcal{T}/\mathcal{C}\to\mathcal{T}'$ making the diagram commute.
The canonical map $\mathcal{Q}_\mathcal{C}$ is called \emph{Verdier localisation} and
the category $\mathcal{T}/\mathcal{C}$ is also written $\mathcal{T}[\mathcal{C}^{-1}]$.
\end{theorem}

% ------------------------------------------------------------------
\subsubsection{Bousfield localisation}

\begin{definition}\label{def:Bousfield-loc}
Let $\mathcal{T}$ be a triangulated category and $\mathcal{S}\subset\mathcal{T}$ a
thick subcategory.  A \emph{Bousfield localisation functor exists for
$\mathcal{S}\subset\mathcal{T}$} if there exists a right adjoint to the natural
functor $F:\mathcal{T}\to\mathcal{T}/\mathcal{S}$.  When this is the case, we call
the right adjoint the \emph{Bousfield localisation functor} and denote it
$G:\mathcal{T}/\mathcal{S}\to\mathcal{T}$.
\end{definition}

\begin{note}\label{note:S-local}
If $\mathcal{S}\subset\mathcal{T}$ is thick and a Bousfield localisation exists, then
by adjunction $\Hom_\mathcal{T}(S,GT)=0$ for every $S\in\mathcal{S}$.
\end{note}

\begin{definition}\label{def:S-local}
An object $T\in\mathcal{T}$ is \emph{$\mathcal{S}$-local} if
$\Hom_\mathcal{T}(S,T)=0$ for every $S\in\mathcal{S}$.
\end{definition}

\begin{proposition}\label{prop:Bousfield-exists}
Let $\mathcal{S}$ be a thick subcategory of a triangulated category $\mathcal{T}$.  A
Bousfield localisation exists for $\mathcal{S}\subset\mathcal{T}$ if and only if the
inclusion $\mathcal{S}\hookrightarrow\mathcal{T}$ admits a right adjoint.
\end{proposition}
\begin{proof}
\cite{Nee01}, Prop.\ 9.1.18.
\end{proof}

\begin{proposition}\label{prop:Bousfield-triangle}
Suppose a Bousfield localisation functor $G:\mathcal{T}/\mathcal{S}\to\mathcal{T}$
exists.  For $T\in\mathcal{T}$ there is a distinguished triangle
\[
  T_\mathcal{S} \longrightarrow T \xrightarrow{\eta_T} GT \longrightarrow T_\mathcal{S}[1]
\]
in which $(1)$ $GT$ is $\mathcal{S}$-local and $(2)$ $T_\mathcal{S}\in\mathcal{S}$.
\end{proposition}

% ------------------------------------------------------------------
\subsection{Neeman's adjoint functor theorem}

\begin{definition}\label{def:compactly-generated}
Let $\mathcal{T}$ be a triangulated category admitting coproducts.
\begin{enumerate}
  \item An object $U\in\mathcal{T}$ is \emph{compact} if
    $\Hom_\mathcal{T}(U,\coprod_\lambda X_\lambda)
    =\coprod_\lambda\Hom_\mathcal{T}(U,X_\lambda)$
    for every coproduct of objects in $\mathcal{T}$.
  \item $\mathcal{T}$ is \emph{compactly generated} if there exists a set
    $\mathcal{U}$ of compact objects such that if
    $\Hom_\mathcal{T}(U,X)=0$ for every $U\in\mathcal{U}$, then $X=0$.
\end{enumerate}
\end{definition}

\begin{definition}\label{def:hocolim}
Let $\mathcal{T}$ be a triangulated category admitting coproducts and
$X_0\to X_1\to X_2\to\cdots$ a sequence in $\mathcal{T}$.  The
\emph{homotopy colimit} $\hocolim X_i$ is defined (up to non-canonical isomorphism)
by the distinguished triangle
\[
  \coprod_i X_i \xrightarrow{\,\id-s\,} \coprod_i X_i \longrightarrow \hocolim X_i
  \longrightarrow \Bigl(\coprod_i X_i\Bigr)[1],
\]
where $s:X_{i_0}\to X_{i_0+1}\to\coprod_i X_i$ is the composition into the factor
$X_{i_0+1}$.
\end{definition}

\begin{proposition}\label{prop:compact-hocolim}
Let $U$ be a compact object in a triangulated category $\mathcal{T}$ admitting
coproducts.  If $X_0\to X_1\to\cdots$ is a sequence in $\mathcal{T}$, then
\[
  \Hom_\mathcal{T}(U,\hocolim X_i) = \colim\Hom_\mathcal{T}(U,X_i).
\]
\end{proposition}
\begin{proof}
\cite{Nee96}, Lem.\ 2.8.
\end{proof}

\begin{theorem}[Brown representability~\cite{Nee96}]\label{thm:Brown}
Let $\mathcal{T}$ be a compactly generated triangulated category and
$H:\mathcal{T}^{op}\to\Ab$ a cohomological functor sending coproducts to products.
Then $H$ is representable.
\end{theorem}

\begin{theorem}[Neeman's adjoint functor theorem~\cite{Nee96}, Thm.\ 4.1]\label{thm:Neeman}
Let $\mathcal{S}$ be a compactly generated triangulated category and $\mathcal{T}$
any triangulated category.  Let $F:\mathcal{S}\to\mathcal{T}$ be a triangulated
functor respecting coproducts.  Then $F$ has a right adjoint
$G:\mathcal{T}\to\mathcal{S}$.
\end{theorem}

% ------------------------------------------------------------------
\subsection{The derived category and derived functors}

\begin{definition}\label{def:complex}
Let $\mathcal{A}$ be an additive category.  A \emph{complex} $A^\bullet$ in
$\mathcal{A}$ is a diagram
\[
  A^\bullet:\quad\cdots\longrightarrow A^{n-1}\xrightarrow{d^{n-1}} A^n
  \xrightarrow{d^n} A^{n+1}\xrightarrow{d^{n+1}}\cdots
\]
with $d^n\circ d^{n-1}=0$ for all $n\in\Z$.  It is \emph{bounded} if
$A^n=0$ for $|n|\gg 0$; \emph{bounded below} if $A^n=0$ for $n\ll 0$;
\emph{bounded above} if $A^n=0$ for $n\gg 0$.
\end{definition}

\begin{definition}\label{def:derived}
Let $\mathcal{A}$ be an abelian category.  Setting
$\mathcal{S}=\{A^\bullet: H^n(A^\bullet)=0\text{ for all }n\in\Z\}$
(the acyclic complexes), $\mathrm{Mor}_\mathcal{S}$ consists of all quasi-isomorphisms
and the \emph{derived category} of $\mathcal{A}$ is
\[
  D(\mathcal{A}) = K(\mathcal{A})[\mathcal{S}^{-1}]
\]
in the sense of Theorem~\ref{thm:Verdier-localisation}.
\end{definition}

\begin{proposition}\label{prop:Ext-Hom}
Let $\mathcal{A}$ be an abelian category with enough injectives.  For $X,Y\in\mathcal{A}$,
\[
  \Hom_{D(\mathcal{A})}(X,Y[i]) =
  \begin{cases}
    \Ext^i_\mathcal{A}(X,Y) & i\geq 0,\\
    0 & i<0.
  \end{cases}
\]
\end{proposition}

% ------------------------------------------------------------------
\subsection{Spectral sequences}
\label{subsec:spectral-sequences}

We follow Boardman's article~\cite{Boa99} closely and supplement it with complete
proofs and diagrams.  The key novelty with respect to classical accounts is the
systematic use of the limit groups $D^\infty$ and $D^{-\infty}$ as the natural
targets for convergence.

% ------------------------------------------------------------------
\subsubsection{Basic definitions}

\begin{definition}\label{def:spectral-sequence}
A \emph{(bigraded cohomological) spectral sequence} is a collection of data:
\begin{enumerate}
  \item abelian groups $E_r^{pq}$ for all $r\geq 1$ and $p,q\in\Z$;
  \item morphisms $d_r^{pq}:E_r^{pq}\to E_r^{p+r,q-r+1}$ with $(d_r)^2=0$;
  \item isomorphisms $E_{r+1}^{pq}\cong H^{pq}(E_r^{**},d_r)
    :=\ker d_r^{pq}\big/\operatorname{im} d_r^{p-r,q+r-1}$.
\end{enumerate}
The integer $n=p+q$ is the \emph{total degree}, $p$ the \emph{filtration degree},
and $q$ the \emph{complementary degree}.
\end{definition}

For each fixed $(p,q)$ we have a nested sequence of subgroups of $E_1^{pq}$:
\[
  0 = B_1^{pq} \subset B_2^{pq} \subset \cdots \subset B_r^{pq}
  \subset \cdots \subset Z_r^{pq} \subset \cdots \subset Z_2^{pq}
  \subset Z_1^{pq} = E_1^{pq},
\]
where $Z_r^{pq}=\ker(d_r^{pq}:E_r^{pq}\to E_r^{p+r,q-r+1})$ and
$B_r^{pq}=\operatorname{im}(d_r^{p-r,q+r-1}:E_r^{p-r,q+r-1}\to E_r^{pq})$
are the \emph{cycles} and \emph{boundaries} at stage $r$, viewed as subgroups
of $E_1^{pq}$ via the successive isomorphisms.

\begin{definition}\label{def:infinite-terms}
The \emph{infinite cycles} and \emph{infinite boundaries} are
\[
  Z_\infty^{pq} = \bigcap_{r\geq 1} Z_r^{pq},\qquad
  B_\infty^{pq} = \bigcup_{r\geq 1} B_r^{pq},
\]
and the \emph{$E_\infty$-term} is the sub-quotient
$E_\infty^{pq} = Z_\infty^{pq}/B_\infty^{pq}$.
\end{definition}

\begin{definition}\label{def:filtration}
Let $H$ be an abelian group.  A \emph{decreasing filtration} of $H$ is a sequence
of subgroups
\[
  \cdots\subset F^{p+1}H\subset F^pH\subset\cdots\subset H,\quad p\in\Z.
\]
Write $F^\infty H:=\bigcap_p F^pH$ and $F^{-\infty}H:=\bigcup_p F^pH$.
The filtration is:
\begin{itemize}
  \item \emph{exhaustive} if $F^{-\infty}H=H$;
  \item \emph{separated} (or Hausdorff) if $F^\infty H=0$;
  \item \emph{complete} if $R^1\!\varprojlim_p F^pH=0$, equivalently
    $H\xrightarrow{\sim}\varprojlim_p H/F^pH$.
\end{itemize}
The \emph{$p$-th graded piece} is $\Gr_F^p H = F^pH/F^{p+1}H$.
\end{definition}

\begin{definition}\label{def:convergence}
Let $E_r^{pq}$ be a spectral sequence and $(H^*,F^*)$ a filtered graded abelian
group.  The spectral sequence \emph{abuts to} (or \emph{converges to})
$(H^*,F^*)$ if there are isomorphisms
\[
  E_\infty^{pq} \xrightarrow{\;\sim\;} \Gr_F^p H^{p+q}
\]
for every $(p,q)$.  One writes $E_1^{pq}\Rightarrow H^n$ (where $n=p+q$).
\begin{enumerate}[label=\textup{(\roman*)}]
  \item The sequence \emph{converges weakly} to $(H^*,F^*)$ if additionally $F^*$
    is exhaustive.
  \item It \emph{converges} if $F^*$ is exhaustive and separated.
  \item It \emph{converges strongly} if $F^*$ is exhaustive, separated, and
    complete.
\end{enumerate}
Strong convergence allows one to recover $H^n$, up to extensions, from the
$E_\infty$-terms via the filtration.
\end{definition}

\begin{definition}\label{def:SS-morphism}
A \emph{morphism of spectral sequences} $f_r^{pq}:E_r^{pq}\to\tilde{E}_r^{pq}$
is a collection of group homomorphisms, one for each $r\geq 1$ and $(p,q)$,
satisfying:
\begin{enumerate}
  \item $f_r^{pq}$ commutes with differentials:
    $\tilde{d}_r^{pq}\circ f_r^{pq}=f_r^{p+r,q-r+1}\circ d_r^{pq}$, i.e.\
    the square
    \[
    \begin{tikzcd}
      E_r^{pq} \ar[r,"d_r^{pq}"] \ar[d,"f_r^{pq}"'] &
        E_r^{p+r,q-r+1} \ar[d,"f_r^{p+r,q-r+1}"] \\
      \tilde{E}_r^{pq} \ar[r,"\tilde{d}_r^{pq}"'] &
        \tilde{E}_r^{p+r,q-r+1}
    \end{tikzcd}
    \]
    commutes for all $r,p,q$.
  \item The induced map on cohomology agrees with $f_{r+1}^{pq}$ under the
    structure isomorphisms $E_{r+1}^{pq}\cong H(E_r^{**},d_r)$.
\end{enumerate}
\end{definition}

\begin{theorem}[Comparison theorem for strongly convergent spectral sequences]
\label{thm:5-term}
Let $E_r^{pq}\Rightarrow H^n$ and $\tilde{E}_r^{pq}\Rightarrow\tilde{H}^n$ be two
spectral sequences converging strongly to the filtered groups $(H^*,F^*)$ and
$(\tilde{H}^*,\tilde{F}^*)$ respectively.  Let
$f_r:E_r^{**}\to\tilde{E}_r^{**}$ be a morphism of spectral sequences and
$f^*:H^*\to\tilde{H}^*$ a morphism of filtered graded groups such that the diagram
\[
\begin{tikzcd}
  E_\infty^{pq} \ar[r,"\sim"] \ar[d,"f_\infty^{pq}"'] &
    \Gr_F^p H^{p+q} \ar[d,"\Gr^p f^{p+q}"] \\
  \tilde{E}_\infty^{pq} \ar[r,"\sim"'] &
    \Gr_{\tilde{F}}^p \tilde{H}^{p+q}
\end{tikzcd}
\]
commutes for all $(p,q)$.  If there exists $a\geq 1$ such that
$f_a^{pq}:E_a^{pq}\xrightarrow{\sim}\tilde{E}_a^{pq}$ is an isomorphism for all
$(p,q)$, then $f_a^{pq}$ is an isomorphism for all $r\geq a$ (in particular at
$r=\infty$), and $f^n:H^n\to\tilde{H}^n$ is an isomorphism of filtered groups for
all $n$.
\end{theorem}
\begin{proof}
Since $f_a^{pq}$ is an isomorphism and commutes with $d_a$, the induced map on
cohomology $f_{a+1}^{pq}:E_{a+1}^{pq}\to\tilde{E}_{a+1}^{pq}$ is also an
isomorphism (the five lemma applied to the short exact sequence
$0\to B_a^{pq}\to Z_a^{pq}\to E_{a+1}^{pq}\to 0$).  By induction, $f_r^{pq}$ is
an isomorphism for all $r\geq a$, hence at $r=\infty$.

The commutative diagram in the statement gives, for each $p$ and $n=p+q$, an
isomorphism $\Gr_F^p H^n\xrightarrow{\sim}\Gr_{\tilde{F}}^p\tilde{H}^n$.  Since
both filtrations are exhaustive, separated, and complete, we apply the following
comparison lemma: a morphism of complete, separated, exhaustive filtrations on $H^n$
and $\tilde{H}^n$ that induces isomorphisms on all graded pieces is itself an
isomorphism of filtered groups.  This follows by induction on the filtration length
(for finite filtrations) and in general by passing to the limit using completeness:
\[
  H^n \cong \varprojlim_p H^n/F^pH^n,\quad
  \tilde{H}^n \cong \varprojlim_p \tilde{H}^n/\tilde{F}^p\tilde{H}^n,
\]
and the induced maps on each finite quotient $H^n/F^pH^n\to
\tilde{H}^n/\tilde{F}^p\tilde{H}^n$ are isomorphisms (by the five lemma applied
to the finite filtration $F^p\supset F^{p+1}\supset\cdots$), so the inverse limit
is an isomorphism.
\end{proof}

% ------------------------------------------------------------------
\subsubsection{Spectral sequences from exact couples}
\label{subsubsec:exact-couples}

\begin{definition}\label{def:exact-couple}
An \emph{exact couple} is a diagram of abelian groups and morphisms
\[
\begin{tikzcd}[column sep=2em, row sep=1.6em]
  D \ar[rr,"\alpha"] && D \ar[dl,"\beta"]\\
  & E \ar[ul,"\gamma"] &
\end{tikzcd}
\]
that is exact at each of the three vertices, i.e.\
$\ker\alpha=\operatorname{im}\gamma$, $\ker\beta=\operatorname{im}\alpha$, and
$\ker\gamma=\operatorname{im}\beta$.
\end{definition}

The composite $d:=\beta\circ\gamma:E\to E$ satisfies $d^2=\beta\gamma\beta\gamma=0$
(since $\gamma\beta=0$ by exactness at $D$).  Thus $(E,d)$ is a differential group.

\begin{proposition}\label{prop:derived-couple}
Given an exact couple $(D,E,\alpha,\beta,\gamma)$, set
\[
  D' := \operatorname{im}\alpha \subset D,\qquad
  E' := H(E,d) = \ker d\,/\,\operatorname{im} d.
\]
Define morphisms $\alpha':D'\to D'$, $\beta':D'\to E'$, $\gamma':E'\to D'$ by
\[
  \alpha'(\alpha x) := \alpha(\alpha x),\quad
  \beta'(\alpha x) := [\beta x],\quad
  \gamma'([e]) := \gamma(e),
\]
where $[{-}]$ denotes the class in $E'$.  Then:
\begin{enumerate}
  \item $\alpha'$, $\beta'$, $\gamma'$ are well-defined.
  \item The derived couple $(D',E',\alpha',\beta',\gamma')$
    \[
    \begin{tikzcd}[column sep=2em, row sep=1.6em]
      D' \ar[rr,"\alpha'"] && D' \ar[dl,"\beta'"]\\
      & E' \ar[ul,"\gamma'"] &
    \end{tikzcd}
    \]
    is again an exact couple.
\end{enumerate}
Iterating, one obtains a sequence of exact couples $(D_r,E_r,\alpha_r,\beta_r,\gamma_r)$
with $D_1=D$, $E_1=E$, and $d_r:=\beta_r\circ\gamma_r:E_r\to E_r$.  The groups
$\{E_r, d_r\}$ form a spectral sequence.
\end{proposition}
\begin{proof}
\textit{Well-definedness of $\beta'$.}  Suppose $\alpha x=\alpha x'$, so
$\alpha(x-x')=0$, hence $x-x'\in\ker\alpha=\operatorname{im}\gamma$, say
$x-x'=\gamma e$.  Then $\beta x - \beta x'=\beta\gamma e = de$, so
$[\beta x]=[\beta x']$ in $E'$.

\textit{Well-definedness of $\gamma'$.}  Suppose $[e]=[e']$, so $e-e'\in\operatorname{im} d
= \operatorname{im}\beta\gamma$.  Write $e-e'=\beta\gamma f$ for some $f\in E$.
Then $\gamma(e)-\gamma(e')=\gamma\beta\gamma f$.  But $\gamma\beta=0$ by exactness,
so $\gamma(e)=\gamma(e')$ and $\gamma'$ is well-defined.  Moreover $\gamma'[e]=
\gamma e\in\ker\alpha=\operatorname{im}\gamma\subset D'$, so $\gamma'$ maps into $D'$.

\textit{Exactness at $D'$ (i.e.\ $\ker\alpha'=\operatorname{im}\gamma'$).}
$\ker\alpha' = \ker(\alpha|_{\operatorname{im}\alpha})
= \operatorname{im}\alpha\cap\ker\alpha
= \operatorname{im}\alpha\cap\operatorname{im}\gamma$.
On the other hand $\operatorname{im}\gamma'
=\gamma(\ker d)=\gamma(\ker\beta\gamma)$.
If $\gamma e\in\operatorname{im}\alpha$ then $\gamma e=\alpha x$ and
$\beta\gamma e=\beta\alpha x=0$ (exactness at $D$), so $e\in\ker\beta\gamma=\ker d$.
Conversely if $e\in\ker d$ then $\beta\gamma e=0$, so $\gamma e\in\ker\beta=
\operatorname{im}\alpha$.  Hence $\ker\alpha'=\operatorname{im}\gamma'$.

\textit{Exactness at $D'$ on the other side
(i.e.\ $\ker\beta'=\operatorname{im}\alpha'$).}
$\ker\beta' = \{\alpha x : [\beta x]=0\} = \{\alpha x : \beta x\in\operatorname{im} d\}
= \{\alpha x : \beta x = \beta\gamma e\text{ for some }e\}$.
Since $\beta(x-\gamma e)=0$, we have $x-\gamma e\in\ker\beta=\operatorname{im}\alpha$,
say $x-\gamma e=\alpha y$.  Then $\alpha x=\alpha^2 y=\alpha'(\alpha y)$.
Hence $\ker\beta'\subset\operatorname{im}\alpha'$; the reverse inclusion is clear.

\textit{Exactness at $E'$.}
$\ker\gamma' = \{[e]\in E': \gamma e=0\} = \{[e]: e\in\ker\gamma\}
= \{[e]: e\in\operatorname{im}\beta\}$.
On the other hand $\operatorname{im}\beta' = \{[\beta x]: x\in D\}
= \beta(\ker\beta\gamma^{-1}(\operatorname{im}\alpha))/\operatorname{im} d$.
Since $\beta\alpha=0$ the image of $\beta$ lands in $\ker\gamma/\operatorname{im} d$
(by exactness at $E$, $\ker\gamma=\operatorname{im}\beta$), so
$\operatorname{im}\beta'=\{[\beta x]\}\subset\ker\gamma'$ and one checks the
reverse inclusion similarly.

Iteration is straightforward: $(D_r,E_r):=(D_{r-1}',E_{r-1}')$ and
$E_{r+1}^{pq}=\ker d_r^{pq}/\operatorname{im}d_r^{p-r,q+r-1}$ recovers the structure
isomorphism of Definition~\ref{def:spectral-sequence}.
\end{proof}

\begin{note}\label{note:bigraded-couple}
In the bigraded case, one indexes the exact couple as
$(D^{**},E^{**},\alpha,\beta,\gamma)$ with
\[
  \deg\alpha = (-1,1),\quad \deg\beta = (0,0),\quad \deg\gamma = (1,0).
\]
After deriving $r-1$ times, the $r$-th couple has degrees
\[
  \deg\alpha_r = (-1,1),\quad
  \deg\beta_r = (r-1,\,1-r),\quad
  \deg\gamma_r = (1,0),
\]
and the differential $d_r=\beta_r\gamma_r$ has bidegree $(r,1-r)$.  In the
homological convention $E_{pq}$ the differential goes
$d_r:E_r^{pq}\to E_r^{p-r,q+r-1}$.
\end{note}

% ------------------------------------------------------------------
\subsubsection{Convergence and limit groups}

Fix a bigraded exact couple $(D^{**},E^{**},\alpha,\beta,\gamma)$ as in
Note~\ref{note:bigraded-couple}.  Write $\alpha:D^{p,*}\to D^{p-1,*}$ for the
component of $\alpha$ in filtration degree.

\begin{definition}\label{def:limit-groups}
Define the \emph{colimit group} and \emph{limit group}:
\[
  D^{-\infty} := \varinjlim_p D^{p,*}
  \quad\text{(under }\alpha\text{)},\qquad
  D^{\infty}  := \varprojlim_p D^{p,*}
  \quad\text{(under }\alpha\text{)}.
\]
Let $\eta^p:D^{p,*}\to D^{-\infty}$ and $\pi^p:D^\infty\to D^{p,*}$ denote the
canonical maps.  Define filtrations
\[
  F^p D^{-\infty} := \operatorname{im}\eta^p,\qquad
  F^p D^{\infty}  := \ker\pi^p.
\]
\end{definition}

\begin{proposition}\label{prop:filtration-properties}
With the notation of Definition~\ref{def:limit-groups}:
\begin{enumerate}
  \item The filtration $F^*D^{-\infty}$ is exhaustive.
  \item The filtration $F^*D^\infty$ is separated and complete.
  \item $\bigcup_p F^p D^\infty = \ker(D^\infty\to D^{-\infty})$.
\end{enumerate}
\end{proposition}
\begin{proof}
(1) follows immediately since $D^{-\infty}=\varinjlim_p D^{p,*}=\bigcup_p\operatorname{im}\eta^p
= \bigcup_p F^pD^{-\infty}$.

(2) \textit{Separated:} $F^\infty D^\infty = \bigcap_p\ker\pi^p$.  An element
$x\in D^\infty$ lies in $\bigcap_p\ker\pi^p$ if and only if $\pi^p(x)=0$ for all
$p$, i.e.\ $x=0$ in the limit, so $F^\infty D^\infty=0$.

\textit{Complete:} We must show $R^1\varprojlim_p F^pD^\infty=0$, i.e.\
$R^1\varprojlim_p\ker\pi^p=0$.  For each $p$ there is a short exact sequence
\[
  0\longrightarrow \ker\pi^p \longrightarrow D^\infty
  \xrightarrow{\pi^p} D^{p,*} \longrightarrow \operatorname{coker}\pi^p\longrightarrow 0.
\]
Taking the inverse limit and using the Milnor sequence for $R^1\varprojlim$, one
checks that $R^1\varprojlim_p D^\infty=0$ (the system $D^\infty\xrightarrow{=}D^\infty
\xrightarrow{=}\cdots$ is Mittag-Leffler), which gives the result.

(3) An element $x\in D^\infty$ satisfies $x\in F^pD^\infty=\ker\pi^p$ for some $p$
if and only if $\pi^p(x)=0\in D^{p,*}$, i.e.\ the image of $x$ under
$D^\infty\to D^{p,*}\to D^{-\infty}$ is zero.  Since the colimit is the union of
the images, $x$ maps to $0$ in $D^{-\infty}$ if and only if it maps to $0$ in
$D^{p,*}$ for some $p$ (as every element of $D^{-\infty}$ comes from some finite
stage), which is precisely $\bigcup_p\ker\pi^p$.
\end{proof}

\begin{remark}\label{rem:RlimD}
There is a short exact sequence of abelian groups
\[
  0\longrightarrow R^1\!\varprojlim_p\, D^{p+1,*}
  \longrightarrow D^\infty
  \longrightarrow \varprojlim_p\, D^{p,*}/D^{p+1,*}
  \longrightarrow 0,
\]
obtained from the Milnor exact sequence applied to the tower $\cdots\to D^{p+1,*}
\to D^{p,*}\to\cdots$.  The term $R^1\!\varprojlim$ is the obstruction to strong
convergence.
\end{remark}

The following key lemma describes the relationship between $E_\infty$ and the
filtration on the limit groups.  It is the technical heart of Boardman's theory.

\begin{lemma}\label{lem:Einf-and-filtration}
There are natural short exact sequences
\begin{align}
  0 &\longrightarrow \operatorname{im}\beta\,/\,B_\infty^{pq}
      \longrightarrow E_\infty^{pq}
      \longrightarrow Z_\infty^{pq}/\operatorname{im}\beta
      \longrightarrow 0, \label{eq:Einf-ses}\\
  0 &\longrightarrow \Gr_F^p D^{-\infty}
      \longrightarrow E_\infty^{pq}
      \longrightarrow RE_\infty^{pq}
      \longrightarrow 0, \label{eq:Einf-RE}
\end{align}
where $RE_\infty^{pq}:=R^1\!\varprojlim_r Z_r^{pq}/B_r^{pq}$ is the derived
$E_\infty$-term.  In particular, if $RE_\infty^{pq}=0$, then
$E_\infty^{pq}\cong\Gr_F^p D^{-\infty}$.
\end{lemma}
\begin{proof}
The Milnor exact sequence applied to the tower
$\cdots\supset Z_r^{pq}\supset Z_{r+1}^{pq}\supset\cdots$ gives
\[
  0\to\varprojlim_r Z_r^{pq}\to Z_1^{pq}\xrightarrow{1-s}
  Z_1^{pq}\to R^1\!\varprojlim_r Z_r^{pq}\to 0,
\]
where $s$ is the shift.  One identifies $\varprojlim_r Z_r^{pq}= Z_\infty^{pq}$
and similarly $\varprojlim_r B_r^{pq}=B_\infty^{pq}$.  The quotient short exact
sequence then yields~\eqref{eq:Einf-ses}.  Sequence~\eqref{eq:Einf-RE} follows from
the identification of $\Gr_F^p D^{-\infty}$ with a quotient of $Z_\infty^{pq}$
via the map $\gamma$; see~\cite{Boa99}, \S5.
\end{proof}

\begin{definition}\label{def:conditional-convergence}
Let $E_r^{pq}$ be the spectral sequence associated to an exact couple
$(D^{**},E^{**},\alpha,\beta,\gamma)$.
\begin{enumerate}
  \item The spectral sequence is \emph{conditionally convergent to $D^\infty$}
    if $D^{-\infty}=0$.
  \item The spectral sequence is \emph{conditionally convergent to $D^{-\infty}$}
    if $D^\infty=0$ and $RD^\infty=0$, where
    $RD^\infty:=R^1\!\varprojlim_p D^{p,*}$.
\end{enumerate}
\end{definition}

\begin{proposition}\label{prop:weak-convergence}
Suppose $E_r^{pq}$ is associated to an exact couple.
\begin{enumerate}
  \item The filtration $F^*D^{-\infty}$ is always exhaustive
    (Proposition~\textup{\ref{prop:filtration-properties}(1)}).
  \item If the spectral sequence converges conditionally to $D^\infty$
    (i.e.\ $D^{-\infty}=0$), then the filtration $F^*D^\infty$ is exhaustive,
    separated, and complete
    (Proposition~\textup{\ref{prop:filtration-properties}(2)}).
  \item In either case the spectral sequence abuts weakly:
    $E_\infty^{pq}\to\Gr_F^pD^{\pm\infty}$ has kernel and cokernel controlled
    by $RE_\infty^{pq}$ and $RD^\infty$ via the exact sequences of
    Lemma~\textup{\ref{lem:Einf-and-filtration}}.
\end{enumerate}
\end{proposition}
\begin{proof}
Part~(1) is immediate from the definition of $D^{-\infty}$ as a colimit.  For
part~(2), when $D^{-\infty}=0$ the exact sequence
$0\to\varprojlim_p F^pD^\infty\to D^\infty\to D^{-\infty}=0$ shows
$\varprojlim_p F^pD^\infty=D^\infty$, i.e.\ the filtration is exhaustive after
reindexing; separatedness and completeness are
Proposition~\ref{prop:filtration-properties}(2).  Part~(3) is a restatement of
Lemma~\ref{lem:Einf-and-filtration}.
\end{proof}

\begin{theorem}[Boardman~\cite{Boa99}, Thm.\ 7.2 and 7.3]\label{thm:convergence-criterion}
Let $E_r^{pq}$ be the spectral sequence of an exact couple.
\begin{enumerate}
  \item \textup{(Strong convergence.)}  If $E_r^{pq}$ converges conditionally to
    $D^\infty$ and $RE_\infty^{pq}=0$ for all $(p,q)$, then $E_r^{pq}$ converges
    strongly to $D^\infty$, i.e.\ there are isomorphisms
    $E_\infty^{pq}\xrightarrow{\sim}\Gr_F^p D^\infty$ and the filtration
    $F^*D^\infty$ is exhaustive, separated, and complete.
  \item \textup{(Weak convergence.)}  If $D^\infty=0$, then $E_r^{pq}$ converges
    weakly to $D^{-\infty}$.
\end{enumerate}
\end{theorem}
\begin{proof}
\textit{Part (1).}  Since $D^{-\infty}=0$, by Proposition~\ref{prop:filtration-properties}
the filtration $F^*D^\infty$ is separated and complete, and the sequence
$0\to D^\infty\to\prod_p D^{p,*}\xrightarrow{\alpha}\prod_p D^{p,*}\to 0$ is
exact (Milnor).  Consider the natural map from the exact couple:
for each $p$ and each $r$, there is a commutative diagram
\[
\begin{tikzcd}[column sep=3em]
  D^{p+r,*} \ar[r,"\alpha^r"] \ar[d] &
    D^{p,*} \ar[r,"\beta"] \ar[d,"\eta^p"] &
    E^{p,*} \ar[r,"\gamma"] \ar[d,dashed,"\phi_r"] &
    D^{p-1,*} \ar[d]\\
  D^{p+r,*} \ar[r,"\alpha^r"'] &
    D^{p,*} \ar[r,"\beta"'] &
    E_r^{p,*} \ar[r,"\gamma_r"'] &
    D_r^{p-1,*}
\end{tikzcd}
\]
where $E_r^{p,*}=\operatorname{im}(\alpha^{r-1}:D^{p+r-1,*}\to D^{p,*})/
\operatorname{im}(\alpha^r:D^{p+r,*}\to D^{p,*})$ and the dashed vertical arrow is
the quotient map.  Passing to the limit $r\to\infty$ yields a map
\[
  \phi: \Gr_F^p D^\infty
  \longrightarrow E_\infty^{pq},
\]
which is an isomorphism when $RE_\infty^{pq}=0$ by the exact sequence
$0\to\Gr_F^p D^{-\infty}\to E_\infty^{pq}\to RE_\infty^{pq}\to 0$
of Lemma~\ref{lem:Einf-and-filtration} together with $D^{-\infty}=0$ (so
$\Gr_F^p D^{-\infty}=\Gr_F^p D^\infty$).

The filtration $F^*D^\infty$ is exhaustive because: for any $x\in D^\infty$, $x$
maps to $0$ in $D^{p,*}$ for all sufficiently small $p$ (since $D^{-\infty}=
\varinjlim_p D^{p,*}=0$ means every element of each $D^{p,*}$ eventually maps to
$0$ under $\alpha$), hence $x\in F^pD^\infty$ for some $p$.

\textit{Part (2).}  The filtration $F^*D^{-\infty}$ is always exhaustive.  When
$D^\infty=0$, separatedness is automatic: $\bigcap_p F^pD^{-\infty}=
\bigcap_p\operatorname{im}\eta^p$.  An element $x$ in this intersection means
$x=\eta^p(x_p)$ for all $p$, with compatible choices; but the compatibility and
$D^\infty=0$ force $x=0$.  The abutment $E_\infty^{pq}\to\Gr_F^pD^{-\infty}$
now follows from Lemma~\ref{lem:Einf-and-filtration} with $RE_\infty^{pq}$
potentially nonzero.
\end{proof}

% ------------------------------------------------------------------
\subsubsection{The spectral sequence associated to a filtered complex}
\label{subsubsec:filtered-complex-SS}

The most classical source of spectral sequences is a filtered complex.

\begin{definition}\label{def:filtered-complex}
Let $\mathcal{A}$ be an abelian category and $K^\bullet$ a cochain complex.
A \emph{decreasing filtration} on $K^\bullet$ is a collection of subcomplexes
\[
  \cdots\supset F^{p+1}K^\bullet\supset F^pK^\bullet\supset\cdots
\]
of $K^\bullet$ (each $F^pK^n$ is a subgroup of $K^n$, and the differential
$d:K^n\to K^{n+1}$ restricts to $d:F^pK^n\to F^pK^{n+1}$).
\end{definition}

\begin{proposition}\label{prop:filtered-complex-SS}
Let $(K^\bullet,F^*)$ be a filtered cochain complex.  There is a spectral sequence
with
\[
  E_0^{pq} = F^pK^{p+q}/F^{p+1}K^{p+q},\qquad
  E_1^{pq} = H^{p+q}(F^pK^\bullet/F^{p+1}K^\bullet),
\]
and $d_r:E_r^{pq}\to E_r^{p+r,q-r+1}$ induced by the differential of $K^\bullet$.
The abutment is $H^*(K^\bullet)$ with the filtration
$F^pH^n(K^\bullet):=\operatorname{im}(H^n(F^pK^\bullet)\to H^n(K^\bullet))$.
\end{proposition}
\begin{proof}
This spectral sequence arises from the exact couple
\[
\begin{tikzcd}[column sep=2em, row sep=1.6em]
  \bigoplus_{p,q}H^{p+q}(F^pK^\bullet) \ar[rr,"\alpha"] &&
  \bigoplus_{p,q}H^{p+q}(F^pK^\bullet) \ar[dl,"\beta"] \\
  & \bigoplus_{p,q}H^{p+q}(F^pK^\bullet/F^{p+1}K^\bullet) \ar[ul,"\gamma"] &
\end{tikzcd}
\]
where $\alpha$ is induced by the inclusions $F^{p+1}K^\bullet\hookrightarrow
F^pK^\bullet$ (shifting the filtration degree by $-1$), $\beta$ is the map on
cohomology induced by the quotient $F^pK^\bullet\to F^pK^\bullet/F^{p+1}K^\bullet$,
and $\gamma$ is the connecting homomorphism from the long exact sequence of
the short exact sequence
\[
  0\longrightarrow F^{p+1}K^\bullet\longrightarrow F^pK^\bullet
  \longrightarrow F^pK^\bullet/F^{p+1}K^\bullet\longrightarrow 0.
\]
Exactness of the couple at each vertex is exactness of this long exact sequence.
One checks that $D^{-\infty}=\varinjlim_p H^n(F^pK^\bullet)=H^n(K^\bullet)$ and
$D^\infty=\varprojlim_p H^n(F^pK^\bullet)$, giving the stated abutment.
\end{proof}

% ------------------------------------------------------------------
\subsubsection{Proposition 1.63: the spectral sequence of a tower}
\label{subsubsec:tower-SS}

The spectral sequence that plays the central role in this thesis arises not from a
filtered complex but from a \emph{tower of objects} in a triangulated category,
via an exact couple built from the long exact sequences of the tower triangles.

Fix a triangulated category $\mathcal{T}$ admitting coproducts, objects
$A,B\in\mathcal{T}$, and a tower
\begin{equation}\label{eq:tower}
  \cdots\longrightarrow A_{n+1}\xrightarrow{f_n} A_n\longrightarrow\cdots
  \longrightarrow A_1\xrightarrow{f_0} A_0
\end{equation}
of objects and morphisms in $\mathcal{T}$.  For each $n$ complete $f_n$ to a
distinguished triangle
\[
  A_{n+1}\xrightarrow{f_n} A_n\longrightarrow Q_n\longrightarrow A_{n+1}[1],
\]
so $Q_n$ is the \emph{$n$-th associated graded object}.

\begin{proposition}\label{prop:tower-SS}
In the above setup, application of $\Hom_{\mathcal{T}}(B,-)$ to the tower
\eqref{eq:tower} yields an exact couple
\[
\begin{tikzcd}[column sep=2.2em, row sep=1.6em]
  \displaystyle\bigoplus_{n,q}\Hom(B,A_n[q])
    \ar[rr,"(f_n)_*"] &&
  \displaystyle\bigoplus_{n,q}\Hom(B,A_n[q])
    \ar[dl,"\beta"] \\
  & \displaystyle\bigoplus_{n,q}\Hom(B,Q_n[q]) \ar[ul,"\gamma"] &
\end{tikzcd}
\]
where $(f_n)_*$ is induced by $f_n$, $\beta$ is induced by $A_n\to Q_n$, and
$\gamma$ is the connecting map $Q_n\to A_{n+1}[1]$.  The associated spectral
sequence has $E_1$-term
\[
  E_1^{p,q} = \Hom_{\mathcal{T}}(B,\,Q_{-p}[q+p])
\]
and converges to $\Hom_{\mathcal{T}}(B,\hocolim_n A_n)$ when $B$ is compact.
\end{proposition}
\begin{proof}
For each $n$, applying the cohomological functor $\Hom_\mathcal{T}(B,-)$
(Lemma~\ref{lem:Hom-homological}) to the distinguished triangle
$A_{n+1}\to A_n\to Q_n\to A_{n+1}[1]$ gives a long exact sequence
\[
\cdots\to\Hom(B,A_{n+1}[q])\xrightarrow{(f_n)_*}\Hom(B,A_n[q])
\xrightarrow{\beta}\Hom(B,Q_n[q])
\xrightarrow{\gamma}\Hom(B,A_{n+1}[q+1])\to\cdots
\]
Assembling these into the bigraded groups
\[
  D^{pq} := \Hom(B,A_{-p}[q+p]),\qquad
  E^{pq} := \Hom(B,Q_{-p}[q+p]),
\]
the maps $\alpha:=(f_n)_*$, $\beta$, $\gamma$ defined entry-wise give the exact
couple, since the long exact sequence provides exactness at each vertex by
definition.  The bidegrees are $\deg\alpha=(-1,1)$, $\deg\beta=(0,0)$,
$\deg\gamma=(1,0)$, matching Note~\ref{note:bigraded-couple}, so the differential
$d_r=\beta_r\gamma_r$ has bidegree $(r,1-r)$.

For abutment: when $B$ is compact, by Proposition~\ref{prop:compact-hocolim},
\[
  \Hom_\mathcal{T}(B,\hocolim_n A_n)
  = \colim_n\Hom_\mathcal{T}(B,A_n)
  = D^{-\infty,*},
\]
where the colimit is taken with respect to the tower maps $(f_n)_*$.  The
filtration on $D^{-\infty}=\Hom(B,\hocolim A_n)$ is then exactly the filtration
$F_\bullet$ from Definition~\ref{def:filtration-F}, and the spectral sequence abuts
to $\Hom(B,\hocolim A_n)$ by Proposition~\ref{prop:weak-convergence}.
\end{proof}

\begin{note}\label{note:tower-halfplane}
If additionally $E_1^{pq}=0$ for all $p$ below some bound $p_0$ (which occurs in
our setting because $B\in DM^{\mathrm{eff}}(s)$ forces
$\Hom(B,bc_{\leq n}A)=0$ for $n\leq s-1$), then the spectral sequence is confined
to a half-plane and the hypothesis $RE_\infty^{pq}=0$ of
Theorem~\ref{thm:convergence-criterion}(1) is automatically satisfied.  This is the
key vanishing used in the proof of Theorem~\ref{thm:main}.
\end{note}

% ------------------------------------------------------------------
\subsection{The Voevodsky category of mixed motives}

Let $k$ be a field.  All schemes are assumed to be separated of finite type over $k$.
We write $\Sch/k$ for the category of schemes over $k$ and $\Sm/k$ for the full
subcategory of smooth schemes.

% ------------------------------------------------------------------
\subsubsection{Finite correspondences}

The canonical reference for this subsection is~\cite{MVW06}, Ch.\ 1.

\begin{definition}\label{def:elementary-correspondence}
Let $X$ be a smooth connected scheme over $k$ and $Y$ any scheme over $k$.  An
\emph{elementary correspondence} from $X$ to $Y$ is an irreducible closed subscheme
$W\subset X\times Y$ which is finite and surjective over $X$.
\end{definition}

The group $\mathrm{Cor}(X,Y)$ is the free abelian group generated by elementary
correspondences from $X$ to $Y$.  Elements of $\mathrm{Cor}(X,Y)$ are called
\emph{finite correspondences}.

\begin{definition}\label{def:Cor}
We define $\Cor:=\Cor(k)$ as the category whose objects are smooth (separated) schemes
of finite type over $k$ and whose morphisms are finite correspondences:
\[
  \Hom_{\Cor}(X,Y) := \Cor(X,Y).
\]
\end{definition}

\begin{definition}\label{def:tensor-Cor}
For $X,Y\in\Cor$, the \emph{tensor product} is $X\otimes Y := X\times_k Y$.  This
makes $\Cor$ a symmetric monoidal category.
\end{definition}

% ------------------------------------------------------------------
\subsubsection{The category of effective geometric motives}

\begin{definition}\label{def:DM-eff-gm}
The category $\DMeffgm$ is defined as the localisation of $K^b(\Cor)$, as a tensor
triangulated category, with respect to:
\begin{itemize}
  \item \textit{Homotopy.}  For $X\in\Sm/k$, one inverts
    $p_*:[X\times\A^1]\to[X]$.
  \item \textit{Mayer–Vietoris.}  For $X=U\cup V$ (Zariski open cover), one inverts
    the canonical map from the cone of $[U\cap V]\to[U]\oplus[V]$ to $[X]$.
\end{itemize}
The category $DM^{\mathrm{eff}}_{\mathrm{gm}}$ of \emph{effective geometric motives} is
the pseudo-abelian completion of $\DMeffgm$ (formally adding kernels and images of
projections).
\end{definition}

\begin{definition}\label{def:DM-gm}
The \emph{category of geometric motives} $\DMgm$ is defined by inverting the functor
$\otimes\Z(1)$ on $DM^{\mathrm{eff}}_{\mathrm{gm}}$; that is, there are objects $X(n)$
for $X\in DM^{\mathrm{eff}}_{\mathrm{gm}}$, $n\in\Z$, with
\[
  \Hom_{\DMgm}(X(n),Y(m)) :=
  \varinjlim_N \Hom_{DM^{\mathrm{eff}}_{\mathrm{gm}}}(X\otimes\Z(N+n),Y\otimes\Z(N+m)).
\]
Here $\Z(1):=\widetilde{[\PP^1]}[-2]$ is the Tate motive and
$\Z(n):=\Z(1)^{\otimes n}$ for $n\geq 0$.
\end{definition}

% ------------------------------------------------------------------
\subsubsection{Sites and sheaves}

\begin{definition}\label{def:presheaf}
A \emph{presheaf} $P$ on a small category $\mathcal{C}$ with values in a category
$\mathcal{A}$ is a functor $P:\mathcal{C}^{op}\to\mathcal{A}$.
\end{definition}

\begin{definition}\label{def:Nisnevich}
Let $X$ be a $k$-variety.  A \emph{Nisnevich cover} $\mathcal{U}\to X$ is an étale
morphism of finite type such that for each separable field extension $F/k$ of finite
degree, the map $\mathcal{U}(F)\to X(F)$ is surjective.
\end{definition}

% ------------------------------------------------------------------
\subsubsection{Motivic complexes and the Suslin complex}

\begin{definition}\label{def:PST}
\begin{enumerate}
  \item The category $\PST:=\PST(k)$ of \emph{presheaves with transfers} is the
    category of additive functors $\Cor^{op}\to\Ab$.
  \item The category $\ShvNis(\Cor)$ of \emph{Nisnevich sheaves with transfers} is
    the full subcategory of $\PST$ consisting of presheaves $F$ such that, for
    each $X\in\Sm/k$, the restriction $F|_{X_{\mathrm{Nis}}}$ is a sheaf.
\end{enumerate}
\end{definition}

\begin{definition}\label{def:Suslin-complex}
For a presheaf $F$ on $\Sm/k$, let $\Delta^n:=\operatorname{Spec} k[t_0,\ldots,t_n]/(\sum t_i -1)$.
We define $C_n(F)(X):=F(X\times\Delta^n)$.  The \emph{Suslin complex} is
\[
  C_*(F)(X):\quad\cdots\to C_1(F)(X)\xrightarrow{d_1} C_0(F)(X),
  \quad d_n=\sum_i(-1)^i\delta_i^*,
\]
where $\delta_i^n:\Delta^{n-1}\to\Delta^n$ is the $i$-th face map.
\end{definition}

\begin{theorem}[Voevodsky]\label{thm:Voevodsky-key}
Let $F\in\PST$ be homotopy invariant over $\Sm/k$.  Then:
\begin{enumerate}
  \item The presheaf $X\mapsto H^q(X_{\mathrm{Nis}},F_{\mathrm{Nis}})$ is a presheaf
    with transfers over $k$.
  \item $F_{\mathrm{Nis}}$ is strictly homotopy invariant.
  \item $F_{\mathrm{Zar}}=F_{\mathrm{Nis}}$ and
    $H^q(X_{\mathrm{Zar}},F_{\mathrm{Zar}})=H^q(X_{\mathrm{Nis}},F_{\mathrm{Nis}})$.
\end{enumerate}
\end{theorem}
\begin{proof}
See~\cite{VSF00}, Ch.\ 3, Thm.\ 4.27 and Thm.\ 5.7, or~\cite{MVW06}, Prop.\ 13.9.
\end{proof}

\begin{theorem}[Voevodsky's embedding theorem]\label{thm:embedding}
There exists a commutative diagram of exact tensor functors
\[
\begin{tikzcd}
  K^b(\Cor) \ar[r,"L"] \ar[d] & D(\ShvNis(\Cor)) \ar[d,"RC_*"]\\
  DM^{\mathrm{eff}}_{\mathrm{gm}} \ar[r,"i"] & \DMeff
\end{tikzcd}
\]
such that $i$ is a full embedding with dense image and $RC_*(L(X))\cong C_*(X)$.
\end{theorem}
\begin{proof}
\cite{VSF00}, Thm.\ 3.2.6.
\end{proof}

% ------------------------------------------------------------------
\subsubsection{Motivic cohomology}

\begin{definition}\label{def:motivic-cohomology}
For $X\in\Sm/k$ and $q\geq 0$, set $M(X)=C_*(X)$ and define
\[
  H^p(X,\Z(q)) := \Hom_{DM^{\mathrm{eff}}_-}\bigl(M_{\mathrm{gm}}(X),\Z(q)[p]\bigr).
\]
\end{definition}

\begin{proposition}\label{prop:motivic-chow}
For $X\in\Sm/k$ and $Y\in\Sch/k$ equidimensional, $\CH^i(Y,j)\cong H^{2i-j}(Y,\Z(i))$
for $i\geq 0$, $j\in\Z$.
\end{proposition}

\begin{proposition}\label{prop:motivic-Z1}
For $X\in\Sm/k$:
\[
  H^n(X,\Z(1)) =
  \begin{cases}
    H^0_{\mathrm{Zar}}(X,\mathcal{O}_X^*) & n=1,\\
    \mathrm{Pic}(X):=H^1_{\mathrm{Zar}}(X,\mathcal{O}_X^*) & n=2,\\
    0 & \text{otherwise}.
  \end{cases}
\]
\end{proposition}

% ------------------------------------------------------------------
\subsubsection{The Voevodsky category of mixed motives}

The triangulated category $DM$ is defined from $\DMeff$ by localising the Tate motive
$\Z(1)\in\DMeff$.  Explicitly:

\begin{definition}\label{def:DM}
There exists a triangulated tensor category $DM$ in which the functor
$\Z(1)\otimes-:DM\to DM$ is an equivalence, together with a full faithful exact
tensor functor $\DMeff\to DM$.
\end{definition}

\begin{note}\label{note:DM-compactly-generated}
By construction (\cite{VSF00}, \S3.2 and~\cite{Ayo07}, Thm.\ 4.5.67), $DM$ is a
compactly generated triangulated category with compact generators
$\{M(X)(n)\}_{n\in\Z}$, where $X$ ranges over smooth schemes of finite type over $k$.
This follows from the fact that the embedding $DM^{\mathrm{eff}}_{\mathrm{gm}}\hookrightarrow DM$
is fully faithful with dense image (Theorem~\ref{thm:embedding}), and the objects
$M(X)(n)$ generate $DM$ under triangles, shifts, and arbitrary coproducts.
\end{note}

\subsubsection{The motivic Postnikov tower}\label{subsubsec:postnikov}

\begin{definition}\label{def:DMeff-n}
Let $DM^{\mathrm{eff}}(n)$ be the localisation subcategory of $\DMeff$ generated by
objects $M(X)(m)[2m]$ with $m\geq n$ and $X\in\Sm/k$.  For $n\in\Z$, we have a
tower of triangulated subcategories
\[
  \cdots\subset DM^{\mathrm{eff}}(n+1)\subset DM^{\mathrm{eff}}(n)\subset\cdots
  \subset DM^{\mathrm{eff}}(0)=\DMeff.
\]
\end{definition}

By Neeman's theorem (Theorem~\ref{thm:Neeman}), the inclusion
$i_n:DM^{\mathrm{eff}}(n)\hookrightarrow DM$ admits a right adjoint $r_n:DM\to
DM^{\mathrm{eff}}(n)$.

\begin{definition}\label{def:slice-tower}
For $E\in\DMeff$, the \emph{slice tower} is
\[
  \cdots\to f_{n+1}E\to f_nE\to f_{n-1}E\to\cdots,\quad f_n:=i_n\circ r_n.
\]
The \emph{$n$-th slice} of $E$ is $s_nE$, defined by the distinguished triangle
\[
  f_{n+1}E\longrightarrow f_nE\longrightarrow s_nE.
\]
\end{definition}

\begin{proposition}\label{prop:slice-Z}
For every $m,n\in\Z$,
\[
  f_m(\Z(n)) =
  \begin{cases}
    \Z(n) & \text{if } m\leq n,\\
    0     & \text{if } m>n.
  \end{cases}
\]
Equivalently, $s_n(\Z(n))=\Z(n)$ and $s_m(\Z(n))=0$ for $m\neq n$.
\end{proposition}
\begin{proof}
Recall from Definition~\ref{def:slice-tower} that $f_m = i_m\circ r_m$, where
$i_m:DM^{\mathrm{eff}}(m)\hookrightarrow DM$ is the inclusion and $r_m:DM\to
DM^{\mathrm{eff}}(m)$ is its right adjoint.  The category $DM^{\mathrm{eff}}(m)$ is
generated by objects $M(X)(j)[2j]$ with $j\geq m$ and $X\in\Sm/k$.

\textit{Case $m\leq n$:} We claim $\Z(n)\in DM^{\mathrm{eff}}(m)$, which gives
$r_m(\Z(n))=\Z(n)$ and hence $f_m(\Z(n))=i_m(\Z(n))=\Z(n)$.  Indeed,
$\Z(n)=\Z(1)^{\otimes n}$ is the $n$-th Tate motive; by construction
(\cite{VSF00}, \S3.1) it belongs to $DM^{\mathrm{eff}}(n)\subset DM^{\mathrm{eff}}(m)$
since $n\geq m$.

More precisely, using Proposition~\ref{prop:bc-properties}(4) adapted to
$DM^{\mathrm{eff}}(m)$: the counit $\epsilon_m^{\Z(n)}:f_m(\Z(n))\to\Z(n)$ is an
isomorphism if and only if $\Z(n)\in DM^{\mathrm{eff}}(m)$.  Since $\Z(n)$ is a
generator of $DM^{\mathrm{eff}}(n)$ and $DM^{\mathrm{eff}}(n)\subset DM^{\mathrm{eff}}(m)$
for $n\geq m$, we conclude $f_m(\Z(n))\cong\Z(n)$.

\textit{Case $m>n$:} We claim $\Hom_{DM}(M(X)(a)[b],\Z(n))=0$ for all $X\in\Sm/k$
and $a\geq m>n$, $b\in\Z$.  By the motivic cohomology calculation
(Proposition~\ref{prop:motivic-Z1} and~\cite{VSF00}, Cor.\ 3.4) together with
the vanishing $H^p(X,\Z(q))=0$ for $q<0$ and for $p>2q$ (\cite{VSF00}, Thm.\ 4.3.7
and~\cite{Voe02b}), one has
\[
  \Hom_{DM}(M(X)(a)[b],\Z(n))
  = \Hom_{DM}(M(X),\Z(n-a)[b])
  = H^b(X,\Z(n-a))
  = 0
\]
since $n-a < 0$ (as $a\geq m>n$).  Therefore $\Z(n)$ is orthogonal to all compact
generators of $DM^{\mathrm{eff}}(m)$, i.e.\ $\Z(n)\in DM^{\mathrm{eff}}(m)^\perp$.

By the definition of $f_m=i_m\circ r_m$ and the adjunction
$\Hom_{DM}(i_mA,E)=\Hom_{DM^{\mathrm{eff}}(m)}(A,r_mE)$, taking $E=\Z(n)$ and
$A\in DM^{\mathrm{eff}}(m)$ arbitrary:
\[
  \Hom_{DM}(A,f_m(\Z(n)))
  = \Hom_{DM}(i_mA,\Z(n))
  = 0,
\]
since $i_mA\in DM^{\mathrm{eff}}(m)$ is built from generators $M(X)(a)[2a]$ with
$a\geq m>n$.  As $DM^{\mathrm{eff}}(m)$ is compactly generated, we conclude
$f_m(\Z(n))=0$.

\textit{Slices:} The distinguished triangle
$f_{n+1}(\Z(n))\to f_n(\Z(n))\to s_n(\Z(n))\to f_{n+1}(\Z(n))[1]$ becomes
$0\to\Z(n)\to s_n(\Z(n))\to 0$ by the two cases, giving $s_n(\Z(n))\cong\Z(n)$.
For $m\neq n$: if $m>n$ then $f_m(\Z(n))=0=f_{m+1}(\Z(n))$, so $s_m(\Z(n))=0$;
if $m<n$ then $f_m(\Z(n))\cong\Z(n)\cong f_{m+1}(\Z(n))$ and the triangle gives
$s_m(\Z(n))=0$.
\end{proof}

% ==================================================================
\section{Results}
% ==================================================================

% ------------------------------------------------------------------
\subsection{The orthogonal tower}
% ------------------------------------------------------------------

Throughout this section we work with motives with coefficients in $R=\Z[\tfrac{1}{p}]$,
where $p$ is the exponential characteristic of the base field $k$.

We recall the construction of the orthogonal tower from~\cite{Pel17}.  By
Note~\ref{note:DM-compactly-generated}, $DM$ is compactly generated.  Let
$\mathcal{G}^{\mathrm{eff}}$ denote the set of compact objects of the form
$\{M(X)(p)\}_{p\geq 0}$ for $X\in\Sm/k$.  For $n\in\Z$, let
$\mathcal{G}^{\mathrm{eff}}(n)$ denote the set of compact objects of the form
$\{M(X)(p)\}_{p\geq n}$ for $X\in\Sm/k$.

Write $DM^\perp(n)$ for the triangulated orthogonal of $DM^{\mathrm{eff}}(n)$
(\cite{Pel17}, Def.\ 2.1.1.):
\[
  DM^\perp(n) = \bigl\{E\in\DMeff : \Hom_{DM^{\mathrm{eff}}}(A,E)=0
    \text{ for all }A\in DM^{\mathrm{eff}}(n)\bigr\}.
\]

By (\cite{Pel17}, Lem.\ 2.1.7(2)), $DM^\perp(n)$ is compactly generated.  By
Theorem~\ref{thm:Neeman}, the inclusion $j_n:DM^\perp(n)\to DM$ admits a right
adjoint
\[
  p_n : DM\longrightarrow DM^\perp(n),
\]
which is triangulated.

\begin{definition}\label{def:orthogonal-cover}
Set $bc_{\leq n}:=j_{n+1}\circ p_{n+1}$ and for $E\in DM$, the \emph{orthogonal
cover} of $E$ is the canonical map $bc_{\leq n}E\to E$.
\end{definition}

\begin{proposition}[\cite{Pel17}, Rem.\ 3.2.5, Prop.\ 3.2.4, Cor.\ 3.2.7]
\label{prop:bc-properties}
Let $E\in DM$.  The orthogonal cover $bc_{\leq n}E$ satisfies:
\begin{enumerate}
  \item $bc_{\leq n}E\in DM^\perp(n+1)$.
  \item The counit $\theta_n:j_{n+1}p_{n+1}\to\id$ is characterised up to unique
    isomorphism by: for all $F\in DM^\perp(n+1)$,
    \[
      \theta_{n*}^E:\Hom_{DM}(F,bc_{\leq n}E)\xrightarrow{\;\sim\;}\Hom_{DM}(F,E).
    \]
  \item There is a canonical natural transformation $\iota_n:bc_{\leq n}\to
    bc_{\leq n+1}$, and $bc_{\leq n}\circ bc_{\leq n+1}\cong bc_{\leq n}$.
  \item For $E\in DM$: $\theta_n^E:bc_{\leq n}E\to E$ is an isomorphism in $DM$
    if and only if $E\in DM^\perp(n+1)$.
\end{enumerate}
\end{proposition}
\begin{proof}
\textit{Part (1).}  By Definition~\ref{def:orthogonal-cover},
$bc_{\leq n}E = j_{n+1}(p_{n+1}E)$.  Since $p_{n+1}:DM\to DM^\perp(n+1)$ maps
into $DM^\perp(n+1)$ by definition, and $j_{n+1}$ is the inclusion
$DM^\perp(n+1)\hookrightarrow DM$, we have $bc_{\leq n}E\in DM^\perp(n+1)$.

\textit{Part (2).}  The unit-counit adjunction $(j_{n+1},p_{n+1})$ gives, for any
$F\in DM$ and $E\in DM$:
\[
  \Hom_{DM}(F, j_{n+1}p_{n+1}E)
  \cong \Hom_{DM^\perp(n+1)}(p_{n+1}F,\, p_{n+1}E).
\]
If $F\in DM^\perp(n+1)$ then the unit $\eta_F:F\to j_{n+1}p_{n+1}F$ is an
isomorphism (since $p_{n+1}$ is right adjoint to the inclusion, and $F$ is already
in $DM^\perp(n+1)$), so
\[
  \Hom_{DM}(F, bc_{\leq n}E)
  = \Hom_{DM}(F, j_{n+1}p_{n+1}E)
  \cong \Hom_{DM}(F, E)
\]
via post-composition with $\theta_n^E:bc_{\leq n}E\to E$ (the counit of the
adjunction at $E$).  This is natural in $F$ and $E$, proving~(2).

\textit{Part (3).}  The containment $DM^\perp(n+2)\subset DM^\perp(n+1)$ (since
$DM^{\mathrm{eff}}(n+1)\subset DM^{\mathrm{eff}}(n+2)$ implies more objects are in the
orthogonal at level $n+1$) gives a natural map between the corresponding projectors.
Concretely, there is a canonical map
\[
\begin{tikzcd}
  bc_{\leq n}E \ar[r,"\iota_n"] \ar[dr,"\theta_n^E"']
    & bc_{\leq n+1}E \ar[d,"\theta_{n+1}^E"]\\
    & E
\end{tikzcd}
\]
defined by the universal property of part~(2): since $bc_{\leq n}E\in DM^\perp(n+1)$
by part~(1), there is a unique lift of $\theta_n^E:bc_{\leq n}E\to E$ through
$\theta_{n+1}^E:bc_{\leq n+1}E\to E$, giving $\iota_n$.  Applying $bc_{\leq n}$ to
the identity $bc_{\leq n+1}E\xrightarrow{\theta_{n+1}^E} E$ and using the universal
property again gives $bc_{\leq n}(bc_{\leq n+1}E)\cong bc_{\leq n}E$.

\textit{Part (4).}  $(\Rightarrow)$:  If $\theta_n^E$ is an isomorphism then
$E\cong bc_{\leq n}E\in DM^\perp(n+1)$ by part~(1).  $(\Leftarrow)$: If
$E\in DM^\perp(n+1)$, then the unit $\eta_E:E\to j_{n+1}p_{n+1}E = bc_{\leq n}E$
is an isomorphism (the adjunction unit is an isomorphism precisely for objects
already in the full subcategory), and $\theta_n^E$ is its inverse.
\end{proof}

\begin{proposition}[\cite{Pel17}, Thm.\ 3.2.12]\label{prop:orthogonal-tower}
For each $E\in DM$, there is a commutative diagram in $DM$:
\[
\begin{tikzcd}[column sep=2.2em, row sep=2.6em]
  \cdots \ar[r]
    & bc_{\leq n}E \ar[r,"\iota_n"] \ar[dr,"\theta_n^E"']
    & bc_{\leq n+1}E \ar[r] \ar[d,"\theta_{n+1}^E"']
    & \cdots \ar[r]
    & \displaystyle\hocolim_{n\to\infty}bc_{\leq n}E \ar[dll,"c"]
  \\
    &
    & E
    &
    &
\end{tikzcd}
\]
where $\iota_n:bc_{\leq n}E\to bc_{\leq n+1}E$ is the canonical natural
transformation of Proposition~\textup{\ref{prop:bc-properties}(3)}, the map
$c:\hocolim_{n\to\infty}bc_{\leq n}E\to E$ is induced by the universal property of
the homotopy colimit applied to the compatible system $\{\theta_n^E\}$, and all
inner triangles commute.
\end{proposition}
\begin{proof}
The existence and properties of the maps $\iota_n$ follow from
Proposition~\ref{prop:bc-properties}(3): there is a canonical natural transformation
$bc_{\leq n}\to bc_{\leq n+1}$ for every $n\in\Z$, and these are compatible in the
sense that the diagram
\[
\begin{tikzcd}
  bc_{\leq n}E \ar[r,"\iota_n"] \ar[dr,"\theta_n^E"']
    & bc_{\leq n+1}E \ar[d,"\theta_{n+1}^E"] \\
    & E
\end{tikzcd}
\]
commutes for each $n$.  Indeed, both $\theta_n^E$ and $\theta_{n+1}^E\circ\iota_n$
are morphisms $bc_{\leq n}E\to E$ in $DM$.  Since $bc_{\leq n}E\in DM^\perp(n+1)$
by Proposition~\ref{prop:bc-properties}(1), the universal property of
$\theta_{n+1}^E$ (Proposition~\ref{prop:bc-properties}(2)) gives a unique factorisation
of any morphism from $bc_{\leq n}E$ to $E$ through $bc_{\leq n+1}E$ — but
$\theta_n^E$ already maps to $E$ directly, and $\iota_n$ followed by $\theta_{n+1}^E$
gives the same map, so both composites agree.

The maps $\theta_n^E:bc_{\leq n}E\to E$ thus form a compatible system indexed over
$\Z$.  By the universal property of the homotopy colimit
(see Definition~\ref{def:hocolim} and~\cite{Nee01}, Prop.\ 1.6.8), there is a
unique morphism
\[
  c:\hocolim_{n\to\infty}bc_{\leq n}E\longrightarrow E
\]
such that for every $n$ the composite
$bc_{\leq n}E\to\hocolim_{n\to\infty}bc_{\leq n}E\xrightarrow{c} E$
equals $\theta_n^E$.  This is precisely the statement that all inner triangles in the
displayed diagram commute.
\end{proof}

\begin{definition}\label{def:orthogonal-tower}
The tower of Proposition~\ref{prop:orthogonal-tower} is called the
\emph{orthogonal tower}.
\end{definition}

\begin{definition}\label{def:filtration-F}
For $E,F\in DM$, consider the increasing filtration $F_\bullet$ on
$\Hom_{DM}(F,E)$ (resp.\ $\Hom_{DM}(F,\hocolim_{n\to\infty}bc_{\leq n}E)$)
where $F_p$ is the image of
$\theta_{p*}^E:\Hom_{DM}(F,bc_{\leq p}E)\to\Hom_{DM}(F,E)$
(resp.\ $\lambda_{p*}^E:\Hom_{DM}(F,bc_{\leq p}E)\to\Hom_{DM}(F,\hocolim bc_{\leq n}E)$).
\end{definition}

% ------------------------------------------------------------------
\subsection{Convergence of the orthogonal tower}
% ------------------------------------------------------------------

\subsubsection{The orthogonal spectral sequence}\label{sec:orthogonal-SS}

Let $E,F\in DM$.  By (\cite{Pel17}, Thm.\ 3.2.8), one constructs canonical
triangulated functors
\[
  bc_{p/p-1}:DM\longrightarrow DM,\quad p\in\Z,
\]
fitting into distinguished triangles in $DM$:
\[
  bc_{\leq p-1}E\longrightarrow bc_{\leq p}E\longrightarrow bc_{p/p-1}E.
\]

\begin{proposition}\label{prop:SS-exists}
The orthogonal tower induces a spectral sequence of homological type
\[
  E^1_{p,q} = \Hom_{DM}\!\bigl(F,(bc_{p/p-1}E)[q-p]\bigr)
  \;\Rightarrow\; \Hom_{DM}(F,E),
\]
with differentials $d_r:E^r_{p,q}\to E^r_{p-r,q-r+1}$, where the abutment carries
the filtration $F_\bullet$ of Definition~\ref{def:filtration-F}.

Similarly, applying $\Hom_{DM}(F,-)$ to the tower $\{bc_{\leq n}E\}$ gives a
spectral sequence converging to $\Hom_{DM}(F,\hocolim_{n\to\infty}bc_{\leq n}E)$:
\begin{equation}\label{eq:SS-hocolim}
  E^1_{p,q} = \Hom_{DM}\!\bigl(F,(bc_{p/p-1}E)[q-p]\bigr)
  \;\Rightarrow\; \Hom_{DM}(F,\hocolim_{n\to\infty}bc_{\leq n}E).
\end{equation}
The canonical map $c:\hocolim bc_{\leq n}E\to E$ of
Proposition~\ref{prop:orthogonal-tower} induces a morphism of spectral sequences
from~\eqref{eq:SS-hocolim} to the first that is the identity on $E_1$-terms.
\end{proposition}
\begin{proof}
Apply the cohomological functor $\Hom_{DM}(F,-)$ to each distinguished triangle
in the orthogonal tower:
\[
  bc_{\leq p-1}E \longrightarrow bc_{\leq p}E \longrightarrow bc_{p/p-1}E
  \longrightarrow (bc_{\leq p-1}E)[1].
\]
Setting $D^{p,q}:=\Hom_{DM}(F,(bc_{\leq -p}E)[q+p])$ and
$E^{p,q}:=\Hom_{DM}(F,(bc_{-p/-p-1}E)[q+p])$, the long exact sequences in
$\Hom_{DM}(F,-)$ assemble into an exact couple
\[
\begin{tikzcd}[column sep=2.8em,row sep=2em]
  \displaystyle\bigoplus_{p,q} D^{p,q}
    \ar[rr,"\alpha"] &&
  \displaystyle\bigoplus_{p,q} D^{p,q}
    \ar[dl,"\beta"] \\[6pt]
  & \displaystyle\bigoplus_{p,q} E^{p,q}
    \ar[ul,"\gamma"] &
\end{tikzcd}
\]
where:
\begin{itemize}
  \item $\alpha:D^{p,q}\to D^{p-1,q+1}$ is induced by the tower map
    $\iota_{-p-1}:bc_{\leq -p-1}E\to bc_{\leq -p}E$;
  \item $\beta:D^{p,q}\to E^{p,q}$ is induced by the quotient map
    $bc_{\leq -p}E\to bc_{-p/-p-1}E$;
  \item $\gamma:E^{p,q}\to D^{p+1,q-1}$ is the connecting map
    $bc_{-p/-p-1}E\to (bc_{\leq -p-1}E)[1]$ from the distinguished triangle.
\end{itemize}
Exactness at each vertex is exactness of the long exact sequence in $\Hom_{DM}(F,-)$
applied to these triangles.  The bidegrees $\deg\alpha=(-1,1)$, $\deg\beta=(0,0)$,
$\deg\gamma=(1,0)$ match Note~\ref{note:bigraded-couple}, so
Proposition~\ref{prop:derived-couple} yields a spectral sequence
$\{E_r^{p,q},d_r\}_{r\geq 1}$ with
\[
  E_1^{p,q} = E^{p,q} = \Hom_{DM}(F,(bc_{p/p-1}E)[q-p]).
\]
The colimit group satisfies $D^{-\infty,q}=\varinjlim_p D^{p,q}
=\Hom_{DM}(F,(\hocolim_n bc_{\leq n}E)[q])$, giving abutment~\eqref{eq:SS-hocolim}.
The morphism $c_*$ is a morphism of exact couples (it commutes with $\alpha$, $\beta$,
$\gamma$ since $c$ commutes with all tower maps and quotients), hence a morphism of
spectral sequences, which is the identity on $E_1$ since $c$ is the identity on each
$bc_{p/p-1}E$ piece.
\end{proof}

\subsubsection{Orthogonality and duality}

Throughout this section, let $Y\in\Sm/k$ be a connected scheme of dimension $d$ and
let $r,s\in\Z$.

\begin{proposition}\label{prop:orthogonality}
$M(Y)(s)[r]\in DM^\perp(d+s+1)$.
\end{proposition}
\begin{proof}
By (\cite{Pel17}, 2.2), it suffices to show that
$\Hom_{DM}(M(X)(a)[b], M(Y)(s)[r])=0$ for all $X\in\Sm/k$ and
$a,b\in\Z$ with $a\geq d+s+1$.  By (\cite{Kel13}, Cor.\ 4.13, Thm.\ 4.12, Thm.\ 5.1)
we may assume the base field $k$ is perfect.  If $k$ admits resolution of
singularities, by (\cite{VSF00}, Thm.\ 4.3.7),
\[
  \Hom_{DM}(M(X)(a)[b],M(Y)(s)[r])
  \cong\Hom_{DM}(M(X)\otimes M^c(Y)(e)[f],\Z)
\]
where $M^c(Y)\in\DMeff$ is the motive of $Y$ with compact support (\cite{VSF00},
\S4.1, Cor.\ 4.1.6), $e=a-s-d$ and $f=b-r-2d$.  For a perfect field of positive
characteristic, the same conclusion follows from (\cite{Sus17}, Thm.\ 5.5.14 and
Lem.\ 5.5.6).  The hypothesis $a\geq d+s+1$ implies $e=a-s-d\geq 1$, so
$M(X)\otimes M^c(Y)(e)[f]\in DM^{\mathrm{eff}}(1)$.

We therefore need to show $\Hom_{DM}(A,\Z)=0$ for $A\in DM^{\mathrm{eff}}(1)$.
Since $DM^{\mathrm{eff}}(1)$ is generated by $\{M(X)(j)[2j]: j\geq 1, X\in\Sm/k\}$,
it suffices to verify this on generators.  Using the motivic cohomology isomorphism
$\Hom_{DM}(M(X)(j)[b],\Z) = H^{b-2j}(X,\Z(-j))$ and the vanishing
$H^p(X,\Z(q))=0$ for $q<0$ (\cite{VSF00}, Cor.\ 3.4), we get zero since $-j\leq -1<0$.
Therefore $M(X)\otimes M^c(Y)(e)[f]\in DM^{\mathrm{eff}}(1)$ maps to zero in
$\Hom_{DM}(-,\Z)$, completing the proof.
\end{proof}

\begin{corollary}\label{cor:bc-isomorphism}
Let $E=M(Y)(s)[r]\in DM$ where $Y\in\Sm/k$ is connected of dimension $d$.  Then:
\begin{enumerate}
  \item The natural map $\theta_{d+s}^E:bc_{\leq d+s}E\to E$ is an isomorphism in
    $DM$.
  \item For any $A\in DM$ and any $f:E\to A$ in $DM$, there exists a unique lifting
    $g:E\to bc_{\leq d+s}A$ such that the diagram
    \[
    \begin{tikzcd}
      & bc_{\leq d+s}A \ar[d,"\theta_{d+s}^A"]\\
      E \ar[ur,"g"] \ar[r,"f"'] & A
    \end{tikzcd}
    \]
    commutes in $DM$.
  \item $f=0$ if and only if $g=0$.
\end{enumerate}
\end{corollary}
\begin{proof}
\textit{Part (1).}  By Proposition~\ref{prop:orthogonality},
$E=M(Y)(s)[r]\in DM^\perp(d+s+1)$.  By Proposition~\ref{prop:bc-properties}(4),
$\theta_{d+s}^E:bc_{\leq d+s}E\to E$ is an isomorphism if and only if
$E\in DM^\perp(d+s+1)$, which holds.

\textit{Part (2).}  Applying part~(1) with $A$ in place of $E$: since
$E=M(Y)(s)[r]\in DM^\perp(d+s+1)$ by Proposition~\ref{prop:orthogonality}, and
$\theta_{d+s}^A:bc_{\leq d+s}A\to A$ satisfies the universal property of
Proposition~\ref{prop:bc-properties}(2), we obtain for any $f:E\to A$ a unique
$g:E\to bc_{\leq d+s}A$ with $\theta_{d+s}^A\circ g = f$.  The diagram commutes
by construction.  Uniqueness follows because two lifts $g,g'$ satisfy
$\theta_{d+s}^A\circ(g-g')=0$, and since $\theta_{d+s*}^A$ is an isomorphism on
$\Hom_{DM}(E,-)$ by Proposition~\ref{prop:bc-properties}(2) (as $E\in DM^\perp(d+s+1)$),
we get $g=g'$.

\textit{Part (3).}  If $f=0$ then by uniqueness in part~(2), $g=0$.  Conversely,
if $g=0$ then $f=\theta_{d+s}^A\circ g=0$.
\end{proof}

\subsubsection{Convergence in $DM$}

In this section we consider objects $A,B\in DM$ where $B=M(X)(s)[r]$ for
$X\in\Sm/k$ and $r,s\in\Z$.

\begin{theorem}\label{thm:DM-strong-convergence}
The spectral sequence
\[
  E^1_{p,q} = \Hom_{DM}\!\bigl(B,(bc_{p/p-1}A)[q-p]\bigr)
  \;\Rightarrow\; \Hom_{DM}(B,\hocolim_{n\to\infty}bc_{\leq n}A)
\]
converges strongly.
\end{theorem}
\begin{proof}
We verify the three conditions of strong convergence
(Definition~\ref{def:convergence}(iii)) for the filtration $F_\bullet$ of
Definition~\ref{def:filtration-F} on $H:=\Hom_{DM}(B,\hocolim_{n\to\infty}bc_{\leq n}A)$.

\textit{Exhaustiveness.}  Since $B=M(X)(s)[r]$ is compact in $DM$
(Note~\ref{note:DM-compactly-generated}), Proposition~\ref{prop:compact-hocolim} gives
\[
  H = \Hom_{DM}(B,\hocolim_{n\to\infty}bc_{\leq n}A)
  \cong \colim_{n\to\infty}\Hom_{DM}(B,bc_{\leq n}A).
\]
Every element of $H$ therefore lies in the image of some
$\Hom_{DM}(B,bc_{\leq p}A)\to H$, i.e.\ in $F_p H$ for some $p$.
Hence $F_{-\infty}H = H$.

\textit{Half-plane vanishing.}  By construction,
$bc_{\leq n}A\in DM^\perp(n+1)$ (Proposition~\ref{prop:bc-properties}(1)).
Since $B=M(X)(s)[r]\in DM^{\mathrm{eff}}(s)$, for $n\leq s-1$ we have
$n+1\leq s$, so $DM^{\mathrm{eff}}(s)\subset DM^{\mathrm{eff}}(n+1)$ and hence
$\Hom_{DM}(B, bc_{\leq n}A)=0$.
Applying $\Hom_{DM}(B,-)$ to the triangle
$bc_{\leq p-1}A\to bc_{\leq p}A\to bc_{p/p-1}A$ therefore gives
$E_1^{p,q}=\Hom_{DM}(B,(bc_{p/p-1}A)[q-p])=0$ for all $p\leq s-1$.
Thus the spectral sequence is confined to the half-plane $\{p\geq s\}$.

\textit{Strong convergence.}  For any fixed total degree $n=p+q$, only finitely
many $p\geq s$ contribute non-zero $E_1^{p,q}$ once we note that
$E_1^{p,q}=\Hom_{DM}(B,(bc_{p/p-1}A)[q-p])$ and the differentials
$d_r:E_r^{p,q}\to E_r^{p-r,q-r+1}$ have source degree $p-r<p$.  Therefore for
$r$ large enough all differentials into and out of position $(p,q)$ are zero, so
$E_\infty^{p,q}=E_r^{p,q}$ for $r\gg 0$, giving $RE_\infty^{p,q}=0$.
By Theorem~\ref{thm:convergence-criterion}(1) (Boardman, applied via
Proposition~\ref{prop:tower-SS} and Note~\ref{note:tower-halfplane}), the spectral
sequence converges strongly.
\end{proof}

\begin{corollary}\label{cor:hocolim-iso}
Suppose the canonical map
\[
  c:\hocolim_{n\to\infty}bc_{\leq n}A\longrightarrow A
\]
induces an isomorphism of abelian groups
\[
  c_*:\Hom_{DM}(B,\hocolim_{n\to\infty}bc_{\leq n}A)\xrightarrow{\;\sim\;}
  \Hom_{DM}(B,A).
\]
Then the spectral sequence
\[
  E^1_{p,q} = \Hom_{DM}(B,(bc_{p/p-1}A)[q-p])\;\Rightarrow\;\Hom_{DM}(B,A)
\]
converges strongly.
\end{corollary}
\begin{proof}
Follows directly from Theorem~\ref{thm:DM-strong-convergence} and the morphism of
spectral sequences~(2) in Section~\ref{sec:orthogonal-SS}.
\end{proof}

The following is the main result of this thesis.

\begin{theorem}\label{thm:main}
The canonical map
\[
  c:\hocolim_{n\to\infty}bc_{\leq n}A\xrightarrow{\;\sim\;} A
\]
is an isomorphism in $DM$.  Consequently, the spectral sequence
\[
  E^1_{p,q} = \Hom_{DM}\!\bigl(B,(bc_{p/p-1}A)[q-p]\bigr)
  \;\Rightarrow\;\Hom_{DM}(B,A)
\]
converges strongly.
\end{theorem}
\begin{proof}
By Corollary~\ref{cor:hocolim-iso}, it suffices to show that $c$ is an isomorphism in
$DM$.  Since $DM$ is compactly generated, it suffices to show that for all $a,b\in\Z$
and all $Y\in\Sm/k$, the map
\[
  c_*:\Hom_{DM}(M(Y)(a)[b],\hocolim_{n\to\infty}bc_{\leq n}A)
  \longrightarrow\Hom_{DM}(M(Y)(a)[b],A)
\]
is an isomorphism of abelian groups.

\textit{Surjectivity.}  Given $f:M(Y)(a)[b]\to A$, by Corollary~\ref{cor:bc-isomorphism}(2)
(with $d=\dim Y$), there is a lifting $g:M(Y)(a)[b]\to bc_{\leq d+a}A$ with
$\theta_{d+a}^A\circ g=f$.  Surjectivity follows from the orthogonal tower
(Proposition~\ref{prop:orthogonal-tower}).

\textit{Injectivity.}  Let $f:M(Y)(a)[b]\to\hocolim_{n\to\infty}bc_{\leq n}A$ with
$c_*(f)=0$.  Since $M(Y)(a)[b]$ is compact, by Proposition~\ref{prop:compact-hocolim},
\[
  \Hom_{DM}(M(Y)(a)[b],\hocolim_{n\to\infty}bc_{\leq n}A)
  \cong\colim_n\Hom_{DM}(M(Y)(a)[b],bc_{\leq n}A),
\]
so $f$ factors as $M(Y)(a)[b]\xrightarrow{f'} bc_{\leq n}A\xrightarrow{\lambda_{n*}^A}
\hocolim bc_{\leq n}A$ for some $n\geq d+a$.  Applying Corollary~\ref{cor:bc-isomorphism}(2)
again, $f'$ factors as
\[
\begin{tikzcd}
  M(Y)(a)[b] \ar[r,"f''"] & bc_{\leq d+a}A\cong bc_{\leq d+a}(bc_{\leq n}A)
    \ar[r] & bc_{\leq n}A.
\end{tikzcd}
\]
By Corollary~\ref{cor:bc-isomorphism}(3), $f''=0$ since
$\theta_{d+a}^A\circ f''= c_*(f)=0$, using Proposition~\ref{prop:bc-properties}(3)
and the two commuting diagrams.  Hence $f=0$.
\end{proof}

% ------------------------------------------------------------------
\subsection{The $\mathbb{A}^1$-stable homotopy category of Morel--Voevodsky}
% ------------------------------------------------------------------

We refer the reader to (\cite{Jar00}, Thm.\ 4.15) for the construction of the
stable model structure on symmetric $T$-spectra.  We write $\SH$ for its homotopy
category, the \emph{$\mathbb{A}^1$-stable homotopy category of Morel--Voevodsky}.

Let $\Sigma_T^\infty X_+\in\SH$, $X\in\Sm/k$, be the infinite suspension of the
simplicial presheaf represented by $X$ with a disjoint basepoint.  By
(\cite{Jar00}, Prop.\ 4.19), $\SH$ is a tensor triangulated category with unit
$\mathbf{1}=\Sigma_T^\infty\operatorname{Spec} k_+$.  Writing $E(1)$ for
$E\otimes\Sigma_T^\infty(\Gm)[-1]$ and $E(n)=(E(n-1))(1)$ inductively for $n\geq 0$,
the functor $E\mapsto E(1)$ is an autoequivalence (\cite{Hov01}, 8.10;
\cite{Ayo07}, Thm.\ 4.3.38).  Like $DM$, $\SH$ is compactly generated with compact
generators
\[
  \mathcal{G}_{\SH} = \bigl\{\Sigma_T^\infty X_+(p): X\in\Sm/k,\; p\in\Z\bigr\}.
\]

\subsubsection{The orthogonal tower in $\SH$}
\label{subsubsec:SH-tower}

The construction of Section~\ref{subsec:spectral-sequences} and §\ref{subsubsec:postnikov}
carries over to $\SH$ verbatim, replacing $DM^{\mathrm{eff}}(n)$ by the analogous
subcategory $\SH^{\mathrm{eff}}(n)$.

\begin{definition}\label{def:SHeff-n}
For $n\in\Z$, let $\SH^{\mathrm{eff}}(n)$ denote the localisation subcategory of $\SH$
generated by $\{\Sigma_T^\infty X_+(p): p\geq n,\; X\in\Sm/k\}$.  Define
\[
  \SH^\perp(n) := \bigl\{E\in\SH : \Hom_{\SH}(A,E)=0
    \text{ for all }A\in\SH^{\mathrm{eff}}(n)\bigr\}.
\]
By (\cite{Jar00}, Thm.\ 4.15 and~\cite{Pel17}, Lem.\ 2.1.7), $\SH^\perp(n)$ is
compactly generated, so by Theorem~\ref{thm:Neeman} the inclusion
$j_n^{\SH}:\SH^\perp(n)\hookrightarrow\SH$ admits a right adjoint
$p_n^{\SH}:\SH\to\SH^\perp(n)$.  Set
\[
  bc_{\leq n}^{\SH} := j_{n+1}^{\SH}\circ p_{n+1}^{\SH}:\SH\longrightarrow\SH,
\]
with counit $\theta_n^{E,\SH}:bc_{\leq n}^{\SH}E\to E$.
\end{definition}

The properties of Proposition~\ref{prop:bc-properties} hold in $\SH$ with identical
proofs (replacing $DM$ by $\SH$ throughout).  In particular, for each $E\in\SH$
there is an orthogonal tower
\[
\begin{tikzcd}[column sep=2.2em, row sep=2.6em]
  \cdots \ar[r]
    & bc_{\leq n}^{\SH}E \ar[r] \ar[dr,"\theta_n^{E,\SH}"']
    & bc_{\leq n+1}^{\SH}E \ar[r] \ar[d,"\theta_{n+1}^{E,\SH}"']
    & \cdots \ar[r]
    & \hocolim_{n\to\infty}bc_{\leq n}^{\SH}E \ar[dll,"c^{\SH}"]
  \\
    &
    & E
    &
    &
\end{tikzcd}
\]
and, for $F,E\in\SH$, a spectral sequence (Proposition~\ref{prop:SS-exists} applied
in $\SH$; see also~\cite{Pel14}):
\[
  E^1_{p,q} = \Hom_{\SH}\!\bigl(F,(bc_{p/p-1}^{\SH}E)[q-p]\bigr)
  \;\Rightarrow\;
  \Hom_{\SH}(F,\hocolim_{n\to\infty}bc_{\leq n}^{\SH}E).
\]
We henceforth drop the superscript $\SH$ and write $bc_{\leq n}$, $\theta_n^E$, etc.\
when the ambient category is clear from context.

\subsubsection{Failure of strong convergence for $\mathbf{1}\in\SH$}

\begin{proposition}\label{prop:counterexample}
The canonical map
\[
  c:\hocolim_{n\to\infty}bc_{\leq n}\mathbf{1}\longrightarrow\mathbf{1}
\]
is \emph{not} an isomorphism in $\SH$.
\end{proposition}
\begin{proof}
Suppose for contradiction that $c$ is an isomorphism.  Since
$\mathbf{1}\in\SH$ is compact (it is among the generators $\mathcal{G}_{\SH}$),
Proposition~\ref{prop:compact-hocolim} gives
\[
  \colim_{n\to\infty}\Hom_{\SH}(\mathbf{1},bc_{\leq n}\mathbf{1})
  \xrightarrow{\;\sim\;}
  \Hom_{\SH}\!\bigl(\mathbf{1},\hocolim_{n\to\infty}bc_{\leq n}\mathbf{1}\bigr)
  \xrightarrow{\;c_*\;}\Hom_{\SH}(\mathbf{1},\mathbf{1}).
\]
Since $c_*$ would be an isomorphism, $\id_{\mathbf{1}}$ is in the image of some
$\Hom_{\SH}(\mathbf{1},bc_{\leq n}\mathbf{1})\to\Hom_{\SH}(\mathbf{1},\mathbf{1})$,
meaning $\id_{\mathbf{1}}$ factors as:
\[
\begin{tikzcd}
  & bc_{\leq n}\mathbf{1} \ar[d,"\theta_n^{\mathbf{1}}"]\\
  \mathbf{1} \ar[ur,"\sigma"] \ar[r,"\id"'] & \mathbf{1}.
\end{tikzcd}
\]
From $\theta_n^{\mathbf{1}}\circ\sigma=\id_{\mathbf{1}}$, we get that $\sigma$ is a
split monomorphism, so $\mathbf{1}$ is a direct summand of $bc_{\leq n}\mathbf{1}$.
Since $bc_{\leq n}\mathbf{1}\in\SH^\perp(n+1)$ by Proposition~\ref{prop:bc-properties}(1)
(in $\SH$), and $\SH^\perp(n+1)$ is closed under direct summands (it is a thick
subcategory by Proposition~\ref{prop:kernel-properties}), we conclude
$\mathbf{1}\in\SH^\perp(n+1)$.

However, by (\cite{Voe02a}, Thm.\ 2.2(3)) $s_m\mathbf{1}=\mathbf{1}_{s_m}\neq 0$
for all $m\geq 0$; more precisely, $\mathbf{1}\in\SH^{\mathrm{eff}}(m)\setminus\SH^{\mathrm{eff}}(m+1)$
for every $m$ (\cite{Lev14}, Conj.\ 9 and p.\ 350), so $\mathbf{1}\notin\SH^\perp(n+1)$
for any $n$.  This is a contradiction.
\end{proof}

\begin{note}\label{note:counterexample-general}
The argument of Proposition~\ref{prop:counterexample} applies to any compact
$A\in\SH$ for which $s_mA\neq 0$ for arbitrarily large $m$.
\end{note}

\begin{corollary}\label{cor:counterexample2}
For the sphere spectrum $\mathbf{1}\in\SH$, the spectral sequence
\[
  E^1_{p,q} = \Hom_{\SH}\!\bigl(\Sigma_T^\infty X_+(s)[t],
    (bc_{p/p-1}\mathbf{1})[q-p]\bigr)
  \;\Rightarrow\;
  \Hom_{\SH}(\Sigma_T^\infty X_+(s)[t],\mathbf{1})
\]
does \emph{not} converge strongly for any $\Sigma_T^\infty X_+(s)[t]$ with
$X\in\Sm/k$, $s,t\in\Z$.
\end{corollary}
\begin{proof}
Suppose for contradiction that the spectral sequence converges strongly for some
generator $B=\Sigma_T^\infty X_+(s)[t]$.  The morphism of spectral sequences induced
by $c_*$ (Proposition~\ref{prop:SS-exists} applied in $\SH$, equation~\eqref{eq:SS-hocolim})
is the identity on $E_1$-terms, so by Theorem~\ref{thm:5-term} the comparison map
\[
  c_*:\Hom_{\SH}(B,\hocolim_{n\to\infty}bc_{\leq n}\mathbf{1})
  \xrightarrow{\;\sim\;} \Hom_{\SH}(B,\mathbf{1})
\]
is an isomorphism of filtered groups.  Since $\SH$ is compactly generated by
$\mathcal{G}_{\SH}$ and this holds for every $B\in\mathcal{G}_{\SH}$
(by the same argument for each generator), $c$ is an isomorphism in $\SH$ — 
contradicting Proposition~\ref{prop:counterexample}.
\end{proof}

\subsubsection{Strong convergence for $\SH^\perp(r)$-objects}

\begin{proposition}\label{prop:SHconvergence}
Let $A\in\SH$ be such that $A\in\SH^\perp(r)$ for some $r\in\Z$.  Then:
\begin{enumerate}
  \item The canonical map $c:\hocolim_{n\to\infty}bc_{\leq n}A\xrightarrow{\;\sim\;} A$
    is an isomorphism in $\SH$.
  \item The spectral sequence
    \[
      E^1_{p,q} = \Hom_{\SH}\!\bigl(\Sigma_T^\infty X_+(s)[t],
        (bc_{p/p-1}A)[q-p]\bigr)
      \;\Rightarrow\;\Hom_{\SH}(\Sigma_T^\infty X_+(s)[t],A)
    \]
    converges strongly for all $X\in\Sm/k$, $s,t\in\Z$.
\end{enumerate}
\end{proposition}
\begin{proof}
\textit{Part (1).}  Let $m\geq r$ be arbitrary.  Since
$DM^{\mathrm{eff}}(m)\supset DM^{\mathrm{eff}}(r)$ is implied by $m\geq r$ (more conditions
generate a larger orthogonal), and $A\in\SH^\perp(r)$ means $\Hom_\SH(B,A)=0$
for all $B\in\SH^{\mathrm{eff}}(r)$, we get $A\in\SH^\perp(m)$ for all $m\geq r$
since $\SH^{\mathrm{eff}}(m)\subset\SH^{\mathrm{eff}}(r)$.

By Proposition~\ref{prop:bc-properties}(4) in $\SH$: since $A\in\SH^\perp(m+1)$
(using $m\geq r$), the counit $\theta_m^A:bc_{\leq m}A\to A$ is an isomorphism.
Therefore in the tower
\[
\begin{tikzcd}[column sep=1.8em]
  \cdots \ar[r] & bc_{\leq r-1}A \ar[r] & bc_{\leq r}A \ar[r]
    \ar[dr,"\sim"'] & bc_{\leq r+1}A \ar[r] \ar[d,"\sim"'] & \cdots \\
  &&& A &&
\end{tikzcd}
\]
all maps $\theta_m^A$ for $m\geq r$ are isomorphisms, and the tower stabilises.
The homotopy colimit of a stable tower is the stable value, so
$\hocolim_{n\to\infty}bc_{\leq n}A\cong A$ and $c$ is an isomorphism.

\textit{Part (2).}  By part~(1) and Corollary~\ref{cor:hocolim-iso} (applied in
$\SH$ with $B=\Sigma_T^\infty X_+(s)[t]$), it suffices that $c_*$ is an isomorphism
on all Hom groups from compact generators.  Since $B$ is compact and $c$ is now an
isomorphism in $\SH$, $c_*$ is immediately an isomorphism.  The strong convergence
then follows from Theorem~\ref{thm:DM-strong-convergence} applied in $\SH$.
\end{proof}

\begin{corollary}\label{cor:sliceconv}
Let $B=\Sigma_T^\infty X_+(s)[t]$ with $X\in\Sm/k$, $s,t\in\Z$.  Consider the
Voevodsky slice spectral sequence (\cite{Voe02a}, \S7) for $G\in\SH$:
\[
  E_1^{m,n} = \Hom_{\SH}(B,s_mG[m+n])\;\Rightarrow\;\Hom_{\SH}(B,G).
\]
The orthogonal spectral sequence
\[
  E^1_{p,q} = \Hom_{\SH}\!\bigl(B,(bc_{p/p-1}s_mG[m+n])[q-p]\bigr)
  \;\Rightarrow\;\Hom_{\SH}(B,s_mG[m+n])
\]
converges strongly towards the $E_1$-term
$E_1^{m,n}=\Hom_{\SH}(B,s_mG[m+n])$.  Moreover, the differential
$d_1:E_1^{m,n}\to E_1^{m+1,n}$ of the slice spectral sequence induces a map
between the orthogonal spectral sequences for consecutive slices, fitting into
the commutative diagram:
\[
\begin{tikzcd}[column sep=3em, row sep=2em]
  E^1_{p,q}\bigl(s_mG[m+n]\bigr) \ar[r,"d_r"] \ar[d,"d_{1*}"'] &
    E^1_{p-r,q-r+1}\bigl(s_mG[m+n]\bigr) \ar[d,"d_{1*}"] \\
  E^1_{p,q}\bigl(s_{m+1}G[m+1+n]\bigr) \ar[r,"d_r"'] &
    E^1_{p-r,q-r+1}\bigl(s_{m+1}G[m+1+n]\bigr).
\end{tikzcd}
\]
\end{corollary}
\begin{proof}
By (\cite{Voe02a}, Thm.\ 2.2(3)), $s_mG\in\SH^\perp(m+1)$ for every $m$ and every
$G\in\SH$.  Shifting: $s_mG[m+n]\in\SH^\perp(m+1)$ as well (shifts preserve
orthogonality since $\SH^\perp(m+1)$ is closed under the translation functor by
the triangulated structure).  Therefore, with $A=s_mG[m+n]$ and $r=m+1$,
Proposition~\ref{prop:SHconvergence}(2) applies and gives strong convergence.

For the morphism of spectral sequences: the differential
$d_1:E_1^{m,n}\to E_1^{m+1,n}$ of the slice spectral sequence is induced by the
natural map $\partial[m+n]:s_mG[m+n]\to s_{m+1}G[m+1+n]$, where $\partial$ is the
composition (\cite{Voe02a}, Thm.\ 2.2(1))
\[
  s_mG \xrightarrow{\sigma_m} f_{m+1}G[1] \xrightarrow{\pi_{m+1}[1]} s_{m+1}G[1].
\]
Since the orthogonal tower functor $bc_{\leq n}^{\SH}$ is triangulated and natural
(Definition~\ref{def:SHeff-n} and Proposition~\ref{prop:bc-properties} in $\SH$),
post-composition with $\partial[m+n]$ defines a morphism of exact couples
(in the sense of Definition~\ref{def:SS-morphism}), hence a morphism of spectral
sequences.  The commutativity of the square with $d_r$ follows from naturality of
the exact couple construction (Proposition~\ref{prop:derived-couple}) with respect
to morphisms of the tower.
\end{proof}



\begin{thebibliography}{99}

\bibitem{SGA4}
\textit{Théorie des topos et cohomologie étale des schémas.}
Tome~1, 2. Séminaire de Géométrie Algébrique du Bois-Marie 1963--1964 (SGA~4).
Dirigé par M.\ Artin, A.\ Grothendieck et J.L.\ Verdier. Lecture Notes in Mathematics
269. Springer-Verlag, Berlin--New York, 1972.

\bibitem{Ayo07}
J.\ Ayoub. \textit{Les six opérations de Grothendieck et le formalisme des cycles
évanescents dans le monde motivique}. II. Astérisque~\textbf{315}:vi+364\,pp.,
2008 (2007).

\bibitem{Bei87}
A.A.\ Beilinson. Height pairing between algebraic cycles. In
\textit{K-theory, Arithmetic and Geometry} (Moscow, 1984--1986), pp.\ 1--25, 1987.

\bibitem{Blo80}
S.\ Bloch. \textit{Lectures on Algebraic Cycles}. Duke University Mathematics Series
IV. Duke University Mathematics Department, Durham, N.C., 1980.

\bibitem{Boa99}
J.M.\ Boardman. Conditionally convergent spectral sequences.
In \textit{Homotopy Invariant Algebraic Structures} (Baltimore, MD, 1998),
Contemp.\ Math.\ \textbf{239}, pp.\ 49--84. Amer.\ Math.\ Soc., Providence, RI,
1999.

\bibitem{BBD}
Y.\ Manin, S.\ Gelfand. \textit{Methods of Homological Algebra}. Springer, 1996.

\bibitem{Mur93}
J.P.\ Murre. On a conjectural filtration on the Chow groups of an algebraic variety.
I. The general conjectures and some examples.
\textit{Indag.\ Math.\ (N.S.)} \textbf{4}(2):177--188, 1993.

\bibitem{Hov01}
M.\ Hovey. Spectra and symmetric spectra in general model categories.
\textit{J.\ Pure Appl.\ Algebra} \textbf{165}(1):63--127, 2001.

\bibitem{Hub06}
A.\ Huber, B.\ Kahn. The slice filtration and mixed Tate motives.
\textit{Compos.\ Math.} \textbf{142}(4):907--936, 2006.

\bibitem{Jar00}
J.F.\ Jardine. Motivic symmetric spectra. \textit{Doc.\ Math.} \textbf{5}:445--553,
2000.

\bibitem{Kel13}
S.\ Kelly. Triangulated categories of motives in positive characteristic (2013).
Preprint, arXiv:1305.5349v2.

\bibitem{MVW06}
C.\ Mazza, V.\ Voevodsky, C.\ Weibel. \textit{Lecture Notes on Motivic Cohomology}.
Clay Mathematics Monographs, vol.\ 2. Amer.\ Math.\ Soc., Providence, RI, 2006.

\bibitem{Lev14}
M.\ Levine. A comparison of motivic and classical stable homotopy theories.
\textit{J.\ Topol.} \textbf{7}(2):327--362, 2014.

\bibitem{Mum69}
D.\ Mumford. Rational equivalence of 0-cycles on surfaces.
\textit{J.\ Math.\ Kyoto Univ.} \textbf{9}:195--204, 1969.

\bibitem{Nee96}
A.\ Neeman. The Grothendieck duality theorem via Bousfield techniques and Brown
representability. \textit{J.\ Amer.\ Math.\ Soc.} \textbf{9}(1):205--236, 1996.

\bibitem{Nee01}
A.\ Neeman. \textit{Triangulated Categories}. Princeton University Press, 2001.

\bibitem{Ner52a}
A.\ Néron. Problèmes arithmétiques et géometriques rattachée à la notion de rang
d'une courbe algébrique dans un corps. \textit{Bull.\ Soc.\ Math.\ France}
\textbf{80}:101--166, 1952.

\bibitem{Ner52b}
A.\ Néron. La théorie de la base pour les diviseurs sur les variétés algébriques.
2nd Coll.\ Géom.\ Alg., Liège, G.\ Thone, pp.\ 119--126, 1952.

\bibitem{Pel14}
P.\ Peláez. \textit{Motivic Birational Covers and Finite Filtrations on Chow
Groups}, 2014.

\bibitem{Pel17}
P.\ Peláez. Mixed motives and motivic birational covers.
\textit{J.\ Pure Appl.\ Algebra} \textbf{221}:1699--1716, 2017.

\bibitem{Sev34}
F.\ Severi. La base per le varietà algebriche di dimensione qualunque contenute in
una data e la teoria generale delle corrispondénze fra i punti di due superficie
algebriche. \textit{Mem.\ Accad.\ Ital.} \textbf{5}:239--283, 1934.

\bibitem{Sus17}
A.\ Suslin. Motivic complexes over nonperfect fields.
\textit{Ann.\ K-Theory} \textbf{2}(2):277--302, 2017.

\bibitem{Stacks}
The Stacks Project. \url{http://stacks.math.columbia.edu} (2017).

\bibitem{Ver77}
J.L.\ Verdier. Catégories triangulées, état~0. Cohomologie étale. Séminaire de
Géométrie Algébrique du Bois-Marie SGA~$4\tfrac{1}{2}$. P.\ Deligne avec la
collaboration de J.F.\ Boutot, A.\ Grothendieck, L.\ Illusie et J.L.\ Verdier.
Lecture Notes in Mathematics~569. Springer-Verlag, 1977.

\bibitem{VSF00}
V.\ Voevodsky. Triangulated categories of motives over a field. In
\textit{Cycles, Transfers, and Motivic Homology Theories}, Ann.\ Math.\ Stud.\
\textbf{143}, pp.\ 188--238. Princeton Univ.\ Press, 2000.

\bibitem{Voe02b}
V.\ Voevodsky. Motivic cohomology groups are isomorphic to higher Chow groups in
any characteristic. \textit{Int.\ Math.\ Res.\ Not.} \textbf{7}:351--355, 2002.

\bibitem{Voe02a}
V.\ Voevodsky. Open problems in the motivic stable homotopy theory. I.
In \textit{Motives, Polylogarithms and Hodge Theory, Part~I} (Irvine, CA, 1998),
Int.\ Press Lect.\ Ser.\ \textbf{3}, pp.\ 3--34. Int.\ Press, Somerville, MA, 2002.

\bibitem{Voe10}
V.\ Voevodsky. Cancellation theorem. \textit{Doc.\ Math.}, Extra Volume: Andrei A.\
Suslin Sixtieth Birthday, pp.\ 671--685, 2010.

\bibitem{VSF00b}
V.\ Voevodsky, A.\ Suslin, E.M.\ Friedlander. \textit{Cycles, Transfers and
Motivic Homology Theories}. Annals of Mathematics Studies \textbf{143}.
Princeton Univ.\ Press, 2000.

\end{thebibliography}

\end{document}
